\subsection{分裂半单李代数的根}

在本节中，\((\mathfrak{g},\mathfrak{h})\) 表示一个分裂半单李代数。

对任意 \(\lambda \in \mathfrak{h}^*\) ，用 \(\mathfrak{g}^{\lambda}(\mathfrak{h})\) ，或简称 \(\mathfrak{g}^{\lambda}\) ，表示 \(\mathfrak{g}\) 相对于 \(\lambda\) 的准素子空间（参见第七章 §1 第 3 节）。回顾 \(\mathfrak{g}^{0} = \mathfrak{h}\) （第七章 §2 第 1 节 命题 4），\(\mathfrak{g}\) 是诸 \(\mathfrak{g}^{\lambda}\) 的直和（第七章 §1 第 3 节 命题 8 与 9），\(\mathfrak{g}^{\lambda}\) 是由 \(x \in \mathfrak{g}\) 中使 \([h, x] = \lambda(h)x\) 对一切 \(h \in \mathfrak{h}\) 成立者组成的集合（第七章 §2 第 4 节 定理 2 之推论 1），以及 \(\mathfrak{h}\) 在 \(\mathfrak{g}\) 上的权是线性形式 \(\lambda\) （\(\mathfrak{h}\) 上的），且 \(\mathfrak{g}^{\lambda} \neq 0\) （第七章 §1 第 1 节）。

定义 2. \((\mathfrak{g}, \mathfrak{h})\) 的根是 h 在 g 上的一个非零权。

用 \(R(\mathfrak{g},\mathfrak{h})\) ，或简称 R，表示 \((\mathfrak{g},\mathfrak{h})\) 的根集。我们有

\[
\mathfrak {g} = \mathfrak {h} \oplus \bigoplus_ {\alpha \in R} \mathfrak {g} ^ {\alpha}.
\]

命题 1. 设 \(\alpha, \beta\) 是 \((\mathfrak{g}, \mathfrak{h})\) 的根，并设 \(\langle \cdot, \cdot \rangle\) 是 \(\mathfrak{g}\) 上一个非退化的不变对称双线性型（例如 \(\mathfrak{g}\) 的 Killing 型）。

\begin{enumerate}
\def\labelenumi{(\roman{enumi})}
\item
  若 \(\alpha + \beta \neq 0\) ，则 \(\mathfrak{g}^{\alpha}\) 与 \(\mathfrak{g}^{\beta}\) 正交。\(\langle \cdot, \cdot \rangle\) 在 \(\mathfrak{g}^{\alpha} \times \mathfrak{g}^{-\alpha}\) 上的限制是非退化的。\(\langle \cdot, \cdot \rangle\) 在 \(\mathfrak{h}\) 上的限制是非退化的。
\item
  设 \(x \in \mathfrak{g}^{\alpha}\) ，\(y \in \mathfrak{g}^{-\alpha}\) ，\(h \in \mathfrak{h}\) 。则 \([x, y] \in \mathfrak{h}\) ，并且
\end{enumerate}

\[
\langle h, [ x, y ] \rangle = \alpha (h) \langle x, y \rangle .
\]

断言 (i) 是第七章 §1 第 3 节 命题 10 (iii) 的一个特例。若 \(x \in \mathfrak{g}^{\alpha}\) ，\(y \in \mathfrak{g}^{-\alpha}\) ，\(h \in \mathfrak{h}\) ，我们有 \([x, y] \in \mathfrak{g}^{\alpha - \alpha} = \mathfrak{h}\) ，并且

\[
\langle h, [ x, y ] \rangle = \langle [ h, x ], y \rangle = \langle \alpha (h) x, y \rangle = \alpha (h) \langle x, y \rangle .
\]

定理 1. 设 \(\alpha\) 是 \((\mathfrak{g},\mathfrak{h})\) 的根。

\begin{enumerate}
\def\labelenumi{(\roman{enumi})}
\item
  向量空间 \(\mathfrak{g}^{\alpha}\) 的维数是 1。
\item
  向量子空间 \(\mathfrak{h}_{\alpha} = [\mathfrak{g}^{\alpha},\mathfrak{g}^{-\alpha}]\) 是 \(\mathfrak{h}\) 的向量子空间，其维数为 1。它包含唯一的元素 \(H_{\alpha}\) 使 \(\alpha (H_{\alpha}) = 2\) 。
\item
  向量子空间 \(\mathfrak{s}_{\alpha} = \mathfrak{h}_{\alpha} + \mathfrak{g}^{\alpha} + \mathfrak{g}^{-\alpha}\) 是 \(\mathfrak{g}\) 的李子代数。
\item
  若 \(X_{\alpha}\) 是 \(\mathfrak{g}^{\alpha}\) 的非零元素，则存在唯一的 \(X_{-\alpha} \in \mathfrak{g}^{-\alpha}\) 使 \([X_{\alpha}, X_{-\alpha}] = -H_{\alpha}\) 。设 \(\varphi\) 是从 \(\mathfrak{sl}(2, k)\) 到 \(\mathfrak{g}\) 的线性映射，它把 \(X_{+}\) 映到 \(X_{\alpha}\) ，把 \(X_{-}\) 映到 \(X_{-\alpha}\) ，把 \(H\) 映到 \(H_{\alpha}\) ；则 \(\varphi\) 是从李代数 \(\mathfrak{sl}(2, k)\) 到李代数 \(\mathfrak{s}_{\alpha}\) 的同构。
\end{enumerate}

\begin{enumerate}
\def\labelenumi{\alph{enumi})}
\item
  设 \(h_{\alpha}\) 是 h 的唯一元素，使 \(\alpha(h) = \langle h_{\alpha}, h \rangle\) 对一切 \(h \in h\) 成立。由命题 1，\([x, y] = \langle x, y \rangle h_{\alpha}\) 对一切 \(x \in g^{\alpha}\) ，\(y \in g^{-\alpha}\) 成立；另一方面 \(\langle g^{\alpha}, g^{-\alpha} \rangle \neq 0\) 。因此 \(h_{\alpha} = [g^{\alpha}, g^{-\alpha}] = kh_{\alpha}\) 。
\item
  选取 \(x \in \mathfrak{g}^{\alpha}\) ，\(y \in \mathfrak{g}^{-\alpha}\) 使 \(\langle x, y \rangle = 1\) ，于是 \([x, y] = h_{\alpha}\) 。回顾 \([h_{\alpha}, x] = \alpha(h_{\alpha})x\) ，\([h_{\alpha}, y] = -\alpha(h_{\alpha})y\) 。若 \(\alpha(h_{\alpha}) = 0\) ，则由此推出 \(kx + ky + kh_{\alpha}\) 是幂零子代数 \(t\) （\(\mathfrak{g}\) 中的）；由于 \(h_{\alpha} \in [t, t]\) ，\(\operatorname{ad}_{\mathfrak{g}} h_{\alpha}\) 是幂零的（第一章 §5 第 3 节 定理 1），而这是荒谬的，因为 \(\operatorname{ad}_{\mathfrak{g}} h_{\alpha}\) 非零且半单。所以 \(\alpha(h_{\alpha}) \neq 0\) 。因此存在唯一的 \(H_{\alpha} \in \mathfrak{h}_{\alpha}\) 使 \(\alpha(H_{\alpha}) = 2\) ，这就证明了 (ii)。
\item
  选取非零元素 \(X_{\alpha}\) （\(\mathfrak{g}^{\alpha}\) 中的）。存在 \(X_{-\alpha} \in \mathfrak{g}^{-\alpha}\) 使 \([X_{\alpha}, X_{-\alpha}] = -H_{\alpha}\) （因为 \([X_{\alpha}, \mathfrak{g}^{-\alpha}] = \mathfrak{h}_{\alpha}\) ，由 \(b\) ）。则
\end{enumerate}

\[
\begin{array}{c} [ H _ {\alpha}, X _ {\alpha} ] = \alpha (H _ {\alpha}) X _ {\alpha} = 2 X _ {\alpha}, [ H _ {\alpha}, X _ {- \alpha} ] = - \alpha (H _ {\alpha}) X _ {- \alpha} = - 2 X _ {- \alpha}, \\ \relax [ X _ {\alpha}, X _ {- \alpha} ] = - H _ {\alpha}; \end{array}
\]

于是 \(kX_{\alpha} + kX_{-\alpha} + kH_{\alpha}\) 是 \(\mathfrak{g}\) 的一个子代数，而线性映射 \(\varphi\) （从 \(\mathfrak{sl}(2,k)\) 到 \(kX_{\alpha} + kX_{-\alpha} + kH_{\alpha}\) ，满足 \(\varphi(X_{+}) = X_{\alpha}, \varphi(X_{-}) = X_{-\alpha}, \varphi(H) = H_{\alpha}\) ）是李代数的一个同构。

\begin{enumerate}
\def\labelenumi{\alph{enumi})}
\setcounter{enumi}{3}
\tightlist
\item
  假设 \(\dim g^{\alpha} > 1\) 。设 y 是 \(g^{-\alpha}\) 的一个非零元素。存在非零元素 \(X_{\alpha}\) （\(g_{\alpha}\) 中的）使 \(\langle y, X_{\alpha} \rangle = 0\) 。按 c) 选取 \(X_{-\alpha}\) ，并考虑表示 \(\rho : u \mapsto \operatorname{ad}_{\mathfrak{g}} \varphi(u)\) （从 \(\mathfrak{sl}(2, k)\) 到 g）。我们有
\end{enumerate}

\[
\begin{array}{c} \rho (H) y = [ \varphi (H), y ] = [ H _ {\alpha}, y ] = - 2 y \\ \rho (X _ {+}) y = [ \varphi (X _ {+}), y ] = [ X _ {\alpha}, y ] = \langle X _ {\alpha}, y \rangle h _ {\alpha} = 0. \end{array}
\]

于是，\(y\) 对 \(\rho\) 是本原元，权为 \(-2\) ，这与 §1 第 2 节 命题 2 矛盾。这就证明了 (i)。

\begin{enumerate}
\def\labelenumi{\alph{enumi})}
\setcounter{enumi}{4}
\tightlist
\item
  断言 (iii) 现在是 c) 的推论。另一方面，若 \(X_{\alpha}\) 是 \(g^{\alpha}\) 的非零元素，则在 c) 中构造的元素 \(X_{-\alpha}\) 是 \(g^{-\alpha}\) 中使 \([X_{\alpha}, X_{-\alpha}] = -H_{\alpha}\) 的唯一元素，因为 \(\dim g^{-\alpha} = 1\) 。(iv) 的最后一个断言是 c) 的推论。证毕。
\end{enumerate}

以下将保留记号 \(h_{\alpha}\) ，\(H_{\alpha}\) ，\(s_{\alpha}\) 。（定义 \(h_{\alpha}\) 时，我们取 \(\langle\cdot,\cdot\rangle\) 等于 Killing 型。）若 \(X_{\alpha}\) 是 \(g^{\alpha}\) 的非零元素，则定理 1 的同构 \(\varphi\) 与表示 \(u \mapsto \operatorname{ad}_{\mathfrak{g}} \varphi(u)\) （\(\mathfrak{sl}(2,k)\) 在 g 上的）称为对应于 \(X_{\alpha}\) 的同构与表示。

推论. 设 \(\Phi\) 是 \(\mathfrak{g}\) 的 Killing 型。对一切 \(a, b \in \mathfrak{h}\) ，

\[
\Phi (a, b) = \sum_ {\gamma \in \mathrm{R}} \gamma (a) \gamma (b).
\]

事实上，ad a.ad b 保持每个 \(g^{\gamma}\) 稳定，且它在 \(g^{\gamma}\) 上的限制是比率为 \(\gamma(a)\gamma(b)\) 的位似；若 \(\gamma\neq0\) ，则 \(\dim g^{\gamma}=1\) 。

命题 2. 设 \(\alpha, \beta \in \mathbb{R}\) 。

\begin{enumerate}
\def\labelenumi{(\roman{enumi})}
\item
  \(\beta(H_{\alpha}) \in \mathbf{Z}\) 。
\item
  若 \(\Phi\) 表示 \(\mathfrak{g}\) 的 Killing 型，则 \(\Phi(H_{\alpha}, H_{\beta}) \in \mathbf{Z}\) 。
\end{enumerate}

设 \(X_{\alpha}\) 是 \(\mathfrak{g}^{\alpha}\) 的非零元素，\(\rho\) 是 \(\mathfrak{sl}(2,k)\) 在 \(\mathfrak{g}\) 上对应于 \(X_{\alpha}\) 的表示。\(\rho(\mathrm{H})\) 的特征值是 0 以及诸 \(\beta(H_{\alpha})\)

其中 \(\beta \in \mathbb{R}\) 。因此 (i) 由 §1 第 2 节 命题 2 的推论得出。断言 (ii) 由 (i) 及定理 1 的推论得出。证毕。

设 \(\alpha \in \mathrm{R}\) ，\(X_{\alpha}\) 是 \(\mathfrak{g}^{\alpha}\) 的非零元素，\(X_{-\alpha}\) 是 \(\mathfrak{g}^{-\alpha}\) 中使 \([X_{\alpha}, X_{-\alpha}] = -H_{\alpha}\) 的元素，\(\rho\) 是 \(\mathfrak{sl}(2, k)\) 在 \(\mathfrak{g}\) 上对应于 \(X_{\alpha}\) 的表示。设 \(\pi\) 是 \(\mathbf{SL}(2, k)\) 在 \(\mathfrak{g}\) 上与 \(\rho\) 相容的表示（ \(\S 1\) 第 4 节 定理 2）。由于 \(\operatorname{ad} X_{\alpha}\) 是幂零的（第七章 \(\S 1\) 第 3 节 命题 10 (iv)），\(\pi(e^{X_+}) = e^{\operatorname{ad} X_{\alpha}}\) 是 \(\mathfrak{g}\) 的初等自同构。类似地，\(\pi(e^{X_-}) = e^{\operatorname{ad} X_{-\alpha}}\) 是 \(\mathfrak{g}\) 的初等自同构。因此 \(\pi(\mathbf{SL}(2, k)) \subset \operatorname{Aut}_e(\mathfrak{g})\) 。我们使用记号 \(\theta(t)\) （\(\S 1\) 第 5 节中的）。对 \(t \in k^*\) ，置

\[
\theta_ {\alpha} (t) = \pi (\theta (t)) = e ^ {\mathrm{ad} t X _ {\alpha}} e ^ {\mathrm{ad} t ^ {- 1} X _ {- \alpha}} e ^ {\mathrm{ad} t X _ {\alpha}}.\tag{1}
\]

引理 1. (i) 对一切 \(h \in \mathfrak{h}\) ，\(\theta_{\alpha}(t).h = h - \alpha(h)H_{\alpha}\) 。

\begin{enumerate}
\def\labelenumi{(\roman{enumi})}
\setcounter{enumi}{1}
\item
  对一切 \(\beta \in \mathrm{R}\) ，\(\theta_{\alpha}(t)(\mathfrak{g}^{\beta}) = \mathfrak{g}^{\beta - \beta (H_{\alpha})\alpha}\) 。
\item
  若 \(\alpha, \beta \in \mathrm{R}\) ，则 \(\beta - \beta(H_{\alpha})\alpha \in \mathrm{R}\) 。
\end{enumerate}

设 \(h \in \mathfrak{h}\) 。若 \(\alpha(h) = 0\) ，则 \([X_{\alpha}, h] = [X_{-\alpha}, h] = 0\) ，于是 \(\theta_{\alpha}(t).h = h\) 。另一方面，§1 第 5 节的公式 (5) 表明 \(\theta_{\alpha}(t).H_{\alpha} = -H_{\alpha}\) 。这就证明了断言 (i)。由此得出 \(\theta_{\alpha}(t)^2 | \mathfrak{h} = \mathrm{Id}\) 。若 \(x \in \mathfrak{g}^{\beta}\) 且 \(h \in \mathfrak{h}\) ，

\[
\begin{array}{r l} [ h, \theta_ {\alpha} (t) x ] & = \theta_ {\alpha} (t). [ \theta_ {\alpha} (t) h, x ] - \beta (\theta_ {\alpha} (t) h). \theta_ {\alpha} (t) x \\ & = (\beta (h) - \alpha (h) \beta (H _ {\alpha})). \theta_ {\alpha} (t) x \\ & = (\beta - \beta (H _ {\alpha}) \alpha) (h). \theta_ {\alpha} (t) x \end{array}
\]

于是 \(\theta_{\alpha}(t)x\in \mathfrak{g}^{\beta -\beta (H_{\alpha})\alpha}\) 。这就证明了 (ii)。断言 (iii) 由 (ii) 得出。

定理 2. (i) 集合 \(\mathrm{R} = \mathrm{R}(\mathfrak{g},\mathfrak{h})\) 是 \(\mathfrak{h}^*\) 中的一个既约根系。

\begin{enumerate}
\def\labelenumi{(\roman{enumi})}
\setcounter{enumi}{1}
\tightlist
\item
  设 \(\alpha \in \mathbb{R}\) 。映射 \(s_{\alpha, H_\alpha}: \lambda \mapsto \lambda - \lambda(H_\alpha)\alpha\) （从 \(\mathfrak{h}^*\) 到 \(\mathfrak{h}^*\) ）是唯一的反射 \(s\) （\(\mathfrak{h}^*\) 中的），使 \(s(\alpha) = -\alpha\) 且 \(s(\mathrm{R}) = \mathrm{R}\) 。对一切 \(t \in k^*\) ，\(s\) 是 \(\theta_\alpha(t)|\mathfrak{h}\) 的转置。
\end{enumerate}

首先，R 生成 \(h^{*}\) ，因为若 \(h \in h\) 使 \(\alpha(h) = 0\) 对一切 \(\alpha \in R\) 成立，则 ad h = 0 ，从而由于 g 的中心为零而得 h = 0 。按定义，\(0 \notin R\) 。设 \(\alpha \in R\) 。由于 \(\alpha(H_{\alpha}) = 2\) ，\(s = s_{\alpha,H_{\alpha}}\) 是使 \(s(\alpha) = -\alpha\) 的一个反射。于是由引理 1 (iii) 有 \(s(R) = R\) ，且 \(\beta(H_{\alpha}) \in \mathbf{Z}\) 对一切 \(\beta \in R\) 成立（命题 2 (i)）。这表明 R 是 \(h^{*}\) 中的一个根系。对一切 \(h \in h\) 与 \(\lambda \in h^{*}\) ，

\[
\langle s (\lambda), h \rangle = \langle \lambda - \lambda (H _ {\alpha}) \alpha , h \rangle = \langle \lambda , h - \alpha (h) H _ {\alpha} \rangle = \langle \lambda , \theta_ {\alpha} (t) h \rangle
\]

于是 \(s\) 是 \(\theta_{\alpha}(t)|\mathfrak{h}\) 的转置。最后，我们证明根系 R 是既约的。设 \(\alpha \in \mathrm{R}\) 且 \(y \in \mathfrak{g}^{2\alpha}\) 。由于 \(3\alpha \notin \mathrm{R}\) （第六章 §1 第 3 节 命题 8），\([X_{\alpha}, y] = 0\) ；另一方面，\([X_{-\alpha}, y] \in \mathfrak{g}^{-\alpha + 2\alpha} = \mathfrak{g}^{\alpha} = kX_{\alpha}\) ，故 \([X_{\alpha}, [X_{-\alpha}, y]] = 0\) ；从而

\[
4 y = 2 \alpha (H _ {\alpha}) y = [ H _ {\alpha}, y ] = - [ [ X _ {\alpha}, X _ {- \alpha} ], y ] = 0
\]

于是 \(y = 0\) 且 \(\mathfrak{g}^{2\alpha} = 0\) 。换言之，\(2\alpha\) 不是根。

证毕。

把 \(\mathfrak{h}\) 与 \(\mathfrak{h}^{**}\) 典则地等同起来。采用第六章 §1 第 1 节的记号，由定理 2 (ii) 我们有

\[
H _ {\alpha} = \alpha^ {\vee} \quad \text { for   all } \alpha \in \mathrm{R}.\tag{2}
\]

诸 \(H_{\alpha}\) 于是构成根系 \(\mathbb{R}^{\vee}\) （\(\mathfrak{h}\) 中的），它与 \(\mathbb{R}\) 互逆。

我们称 \(\mathrm{R}(\mathfrak{g},\mathfrak{h})\) 为 \((\mathfrak{g},\mathfrak{h})\) 的根系。反射 \(s_{\alpha,H_{\alpha}}\) 将简记为 \(s_{\alpha}\) 。Weyl 群、权群、Coxeter 数 \ldots{} （\(\mathrm{R}(\mathfrak{g},\mathfrak{h})\) 的）称为 Weyl 群、权群、Coxeter 数 \ldots{} （\((\mathfrak{g},\mathfrak{h})\) 的）。如同第六章 §1 第 1 节，我们把 Weyl 群看作不仅作用于 \(h^{*}\) ，而且由结构转移也作用于 h ，使得 \(s_{\alpha} = \theta_{\alpha}(t)|\mathfrak{h}\) 。由于诸 \(\theta_{\alpha}(t)\) 是 g 的初等自同构，我们有：

推论. \((\mathfrak{g},\mathfrak{h})\) 的 Weyl 群的每个元素在 h 上作用时，都是 g 的某个初等自同构在 h 上的限制。

关于此结果的一个逆命题，见 §5 第 2 节 命题 4。

注 1. 若用 \(\mathfrak{h}_{\mathbf{Q}}\) （相应地 \(\mathfrak{h}_{\mathbf{Q}}^{*}\) ）表示 \(\mathbf{Q}\) -向量空间 \(\mathfrak{h}\) （相应地 \(\mathfrak{h}^{*}\) ）中由诸 \(H_{\alpha}\) （相应地诸 \(\alpha\) ，其中 \(\alpha \in \mathbb{R}\) ）生成的子空间，则 \(\mathfrak{h}\) （相应地 \(\mathfrak{h}^{*}\) ）可以典则地等同于 \(\mathfrak{h}_{\mathbf{Q}} \otimes_{\mathbf{Q}} k\) （相应地 \(\mathfrak{h}_{\mathbf{Q}}^{*} \otimes_{\mathbf{Q}} k\) ），并且 \(\mathfrak{h}_{\mathbf{Q}}^{*}\) 可以等同于 \(\mathfrak{h}_{\mathbf{Q}}\) 的对偶（第六章 §1 第 1 节 命题 1）。我们称 \(\mathfrak{h}_{\mathbf{Q}}\) 与 \(\mathfrak{h}_{\mathbf{Q}}^{*}\) 为 \(\mathbf{Q}\) 在 \(\mathfrak{h}\) 与 \(\mathfrak{h}^{*}\) 上的典则结构（代数 第二章 §8 第 1 节 定义 1）。当提到 \(\mathbf{Q}\) -有理性（对 \(\mathfrak{h}\) 的向量子空间、对 \(\mathfrak{h}\) 上的线性形式等等）时，除非另有说明，我们指的都是这些结构。当提到 R(g, h) 的 Weyl 房或面时，我们将在 \(\mathfrak{h}_{\mathbf{Q}} \otimes_{\mathbf{Q}} \mathbf{R}\) 或 \(\mathfrak{h}_{\mathbf{Q}}^{*} \otimes_{\mathbf{Q}} \mathbf{R}\) 中工作，并把它们记作 \(\mathfrak{h}_{\mathbf{R}}\) 与 \(\mathfrak{h}_{\mathbf{R}}^{*}\) 。

注 2. 根系 \(\mathbb{R}^{\vee}\) （\(\mathfrak{h}\) 中的）定义了一个非退化对称双线性型 \(\beta\) （\(\mathfrak{h}\) 上的）（第六章 §1 第 1 节 命题 3），即形式 \((a,b)\mapsto \sum_{\alpha \in \mathbb{R}}\langle \alpha ,a\rangle \langle \alpha ,b\rangle\) 。由定理 1 的推论，这个形式正是 Killing 型在 \(\mathfrak{h}\) 上的限制。\(\beta |\mathfrak{h}_{\mathbf{Q}}\times \mathfrak{h}_{\mathbf{Q}}\) 到 \(\mathfrak{h}_{\mathbf{Q}}\otimes_{\mathbf{Q}}\mathbf{R}\) 上的扩张是正定的、非退化的（第六章 §1 第 1 节 命题 3）。另一方面，我们可以看到，\(\mathfrak{h}^*\) 上的逆形式（对应于 \(\mathfrak{h}\) 上的一个限制，即 \(\mathfrak{g}\) 的 Killing 型的限制）就是 R 的典则双线性型 \(\Phi_{\mathrm{R}}\) （第六章 §1 第 12 节）。

设 \((\mathfrak{g}_{1},\mathfrak{h}_{1})\) ， \((\mathfrak{g}_{2},\mathfrak{h}_{2})\) 是分裂半单李代数，\(\varphi\) 是从 \(\mathfrak{g}_{1}\) 到 \(\mathfrak{g}_{2}\) 的同构，且 \(\varphi(\mathfrak{h}_{1})=\mathfrak{h}_{2}\) 。由结构转移，映射 \(\varphi|\mathfrak{h}_{1}\) 的转置把 \(R(\mathfrak{g}_{2},\mathfrak{h}_{2})\) 变到 \(R(\mathfrak{g}_{1},\mathfrak{h}_{1})\) 。

命题 3. 设 g 是半单李代数，\(h_{1}\) 与 \(h_{2}\) 是 g 的分裂嘉当子代数。存在从 \(h_{1}^{*}\) 到 \(h_{2}^{*}\) 的同构，它把 \(\mathrm{R}(\mathfrak{g},\mathfrak{h}_{1})\) 变到 \(\mathrm{R}(\mathfrak{g},\mathfrak{h}_{2})\) 。

（更精确的结果，见 §3 第 3 节 命题 10 的推论，以及 §5 第 3 节 命题 5）。

设 \(k'\) 是 \(k\) 的一个代数闭包，\(\mathfrak{g}' = \mathfrak{g} \otimes_k k'\) ，\(\mathfrak{h}_i' = \mathfrak{h}_i \otimes_k k'\) 。那么 \(\mathrm{R}(\mathfrak{g}', \mathfrak{h}_i')\) 是 \(\mathrm{R}(\mathfrak{g}, \mathfrak{h}_i)\) 在映射 \(\lambda \mapsto \lambda \otimes 1\) （从 \(\mathfrak{h}_i^*\) 到 \(\mathfrak{h}_i^* \otimes_k k' = \mathfrak{h}_i'^*\) ）下的像。由第七章 §3 第 2 节 定理 1，存在 \(\mathfrak{g}'\) 的自同构把 \(\mathfrak{h}_1'\) 变到 \(\mathfrak{h}_2'\) ，因而存在同构 \(\varphi\) （从 \(\mathfrak{h}_1'^*\) 到 \(\mathfrak{h}_2'^*\) ），它把 \(\mathrm{R}(\mathfrak{g}', \mathfrak{h}_1')\) 变到 \(\mathrm{R}(\mathfrak{g}', \mathfrak{h}_2')\) 。于是 \(\varphi | \mathfrak{h}_1^*\) 把 \(\mathrm{R}(\mathfrak{g}, \mathfrak{h}_1)\) 变到 \(\mathrm{R}(\mathfrak{g}, \mathfrak{h}_2)\) ，从而把 \(\mathfrak{h}_1^*\) 变到 \(\mathfrak{h}_2^*\) 。证毕。

由命题 3，\((\mathfrak{g},\mathfrak{h})\) 的根系在同构意义下只依赖于 g 而不依赖于 h。同样，Weyl 群、权群 \ldots{} （\((\mathfrak{g},\mathfrak{h})\) 的）按语言滥用简称为 g 的 Weyl 群、权群 \ldots{} （另参见 §5 第 3 节 注 2）。若 g 的邓金图是 \(A_{l}\) 型，或 \(B_{l}\) 型，\ldots{} （参见第六章 §4 第 2 节 定理 3），我们就说 g 是 \(A_{l}\) 型，或 \(B_{l}\) 型，\ldots 型的。

回顾，若 \(\alpha\) 与 \(\beta\) 是线性无关的根，则 \(j \in \mathbf{Z}\) 中使 \(\beta + j\alpha \in \mathrm{R}\) 者的集合是区间 \([-q, p]\) （\(\mathbf{Z}\) 中的，包含 0），且 \(p - q = -\langle \beta, \alpha^{\vee} \rangle = -\beta(H_{\alpha})\) （第六章 §1 第 3 节 命题 9）。

命题 4. 设 \(\alpha\) 与 \(\beta\) 是线性无关的根。设 \(p\) （相应地 \(q\) ）是最大的整数 \(j\) ，它使 \(\beta + j\alpha\) （相应地 \(\beta - j\alpha\) ）为根。

\begin{enumerate}
\def\labelenumi{(\roman{enumi})}
\item
  向量子空间 \(\sum_{-q\leq j\leq p}\mathfrak{g}^{\beta +j\alpha}\) （\(\mathfrak{g}\) 中的）是单 \(\mathfrak{s}_{\alpha}\) -模，维数为 \(p + q + 1\) 。
\item
  若 \(\alpha + \beta\) 是根，则 \([\mathfrak{g}^{\alpha}, \mathfrak{g}^{\beta}] = \mathfrak{g}^{\alpha + \beta}\) 。
\end{enumerate}

设 \(X_{\alpha}\) （相应地 x）是 \(\mathfrak{g}^{\alpha}\) （相应地 \(\mathfrak{g}^{\beta+p\alpha}\) ）的非零元素。则

\[
\begin{array}{c} {[ X _ {\alpha}, x ] \in \mathfrak {g} ^ {\beta + (p + 1) \alpha} = 0} \\ {[ H _ {\alpha}, x ] = (\beta (H _ {\alpha}) + p \alpha (H _ {\alpha})) x = (- p + q + 2 p) x = (p + q) x.} \end{array}
\]

于是，\(x\) 是权为 \(p + q\) 的本原元（对 \(\mathfrak{sl}(2,k)\) 在 \(\mathfrak{g}\) 上对应于 \(X_{\alpha}\) 的表示而言）；而 \(\mathfrak{sl}(2,k)\) -模 \(\sum_{-q\leq j\leq p}\mathfrak{g}^{\beta +j\alpha}\) 的维数为 \(p + q + 1\) ；因此它是单模（§1 第 2 节 命题 2）。若 \(\alpha +\beta \in \mathrm{R}\) ，则 \(p\geq 1\) ，于是 \(\mathfrak{g}^{\beta}\) 的元素不是本原的，从而 \([X_{\alpha},\mathfrak{g}^{\beta}]\neq 0\) 。由于 \([\mathfrak{g}^{\alpha},\mathfrak{g}^{\beta}]\subset \mathfrak{g}^{\alpha +\beta}\) ，我们最终得到 \([\mathfrak{g}^{\alpha},\mathfrak{g}^{\beta}] = \mathfrak{g}^{\alpha +\beta}\) 。

注 3. 回顾，由第六章 §1 第 3 节 命题 9 的推论，整数 \(p + q + 1\) 只能取值 1, 2, 3, 4。

注 4. 设 \((\mathfrak{g},\mathfrak{h})\) 是分裂约化李代数，\(\mathfrak{c}\) 是 \(\mathfrak{g},\mathfrak{g}' = \mathcal{D}\mathfrak{g}\) 中 g 的中心，\(\mathfrak{h}' = \mathfrak{h} \cap \mathfrak{g}'\) 。则 \(\mathfrak{h} = \mathfrak{c} \times \mathfrak{h}'\) ，并且我们把 \(\mathfrak{h}'^*\) 等同于 \(\mathfrak{h}^*\) 的一个向量子空间。对任意 \(\lambda \in \mathfrak{h}^*\) ，\(\lambda \neq 0\) ，准素子空间 \(\mathfrak{g}^\lambda\) （相对于 \(\lambda\) 的）等于 \(\mathfrak{g}'^{\lambda | \mathfrak{h}'}\) 。\(\mathfrak{h}\) 在 \(\mathfrak{g}\) 上的非零权称为 \((\mathfrak{g},\mathfrak{h})\) 的根；每个根在 \(\mathfrak{c}\) 上为零。用 \(\mathrm{R}(\mathfrak{g},\mathfrak{h})\) 表示 \((\mathfrak{g},\mathfrak{h})\) 的根集；它可以典则地等同于 \(\mathrm{R}(\mathfrak{g}',\mathfrak{h}')\) 。设 \(\alpha \in \mathrm{R}(\mathfrak{g},\mathfrak{h})\) 。我们如同半单情形那样定义 \(h_\alpha\) ，\(H_\alpha\) ，\(\mathfrak{s}_\alpha\) ，同构 \(\mathfrak{sl}(2,k) \to \mathfrak{s}_\alpha\) ，以及 \(\mathfrak{sl}(2,k)\) 在 \(\mathfrak{g}\) 上对应于 \(\alpha\) 的表示。

\subsection{不变双线性型}

命题 5. 设 \((\mathfrak{g},\mathfrak{h})\) 是分裂半单李代数，\(\Phi\) 是 g 上的不变对称双线性型，W 是 \((\mathfrak{g},\mathfrak{h})\) 的 Weyl 群。则 \(\Phi'\) （\(\Phi\) 在 h 上的限制）在 W 下不变。而且，若 \(\Phi\) 非退化，则 \(\Phi'\) 也非退化。

设 \(\alpha \in \mathbb{R}\) ，设 \(X_{\alpha}\) 是 \(\mathfrak{g}^{\alpha}\) 的非零元素，\(\rho\) 是对应的 \(\mathfrak{sl}(2,k)\) 在 \(\mathfrak{g}\) 上的表示，\(\pi\) 是 \(\mathbf{SL}(2,k)\) 在 \(\mathfrak{g}\) 上与 \(\rho\) 相容的表示。则 \(\varPhi\) 在 \(\rho\) 下不变，从而在 \(\pi\) 下不变（§1 第 4 节）。特别地，\(\varPhi'\) 在 \(\theta_{\alpha}(t)|\mathfrak{h}\) （第 2 节）下不变，从而在 W 下不变。最后一个断言由命题 1 (i) 得出。

命题 6. 设 \((\mathfrak{g},\mathfrak{h})\) 是分裂半单李代数，\(\varPhi\) 是 \(\mathfrak{g}\) 上的非退化不变对称双线性型。对一切 \(\alpha \in \mathbb{R}\) ，设 \(X_{\alpha}\) 是 \(\mathfrak{g}^{\alpha}\) 的非零元素。设 \((H_i)_{i\in I}\) 是 \(\mathfrak{h}\) 的基，\((H_i')_{i\in I}\) 是 \(\mathfrak{h}\) 中使 \(\varPhi(H_i,H_j') = \delta_{ij}\) 的基。这时，对应于 \(\varPhi\) 的 Casimir 元（\(\mathfrak{g}\) 的包络代数中的，第一章 §3 第 7 节）是

\[
\sum_ {\alpha \in \mathrm{R}} \frac {1}{\varPhi (X _ {\alpha} , X _ {- \alpha})} X _ {\alpha} X _ {- \alpha} + \sum_ {i \in I} H _ {i} H _ {i} ^ {\prime}.
\]

事实上，由命题 1，\(\Phi(H_i, X_\alpha) = \Phi(H_i', X_\alpha) = 0\) 对一切 \(i \in I\) ，\(\alpha \in \mathrm{R}\) 成立，并且 \(\Phi\left(\frac{1}{\Phi(X_\alpha, X_{-\alpha})} X_\alpha, X_{-\beta}\right) = \delta_{\alpha\beta}\) 对一切 \(\alpha, \beta \in \mathrm{R}\) 成立。

\subsection{\texorpdfstring{系数 \(\mathbf{N}_{\alpha \beta}\)}{系数 \textbackslash mathbf\{N\}\_\{\textbackslash alpha \textbackslash beta\}}}

在本节中，我们仍用 \((\mathfrak{g},\mathfrak{h})\) 表示一个分裂半单李代数。

引理 2. 存在一个族 \((X_{\alpha})_{\alpha \in \mathbb{R}}\) ，使对一切 \(\alpha \in \mathbb{R}\) ，

\[
X _ {\alpha} \in \mathfrak {g} ^ {\alpha} \text { and } [ X _ {\alpha}, X _ {- \alpha} ] = - H _ {\alpha}.
\]

设 \(R_{1}\) 是 R 的子集，使 \(R = R_{1} \cup (-R_{1})\) 且 \(R_{1} \cap (-R_{1}) = \varnothing\) 。对 \(\alpha \in R_{1}\) ，任取非零元素 \(X_{\alpha}\) （\(g^{\alpha}\) 中的）。存在唯一的 \(X_{-\alpha} \in g^{-\alpha}\) 使 \([X_{\alpha}, X_{-\alpha}] = -H_{\alpha}\) （定理 1 (iv)）。则

\[
\left[ X _ {- \alpha}, X _ {\alpha} \right] = H _ {\alpha} = - H _ {- \alpha}.\tag{Q.E.D.}
\]

若 \((X_{\alpha})_{\alpha \in \mathbb{R}}\) 是满足引理 2 的条件的一个族，则满足这些条件的最一般的族是 \((t_{\alpha}X_{\alpha})_{\alpha \in \mathbb{R}}\) ，其中 \(t_{\alpha}\in k^{*}\) 且 \(t_{\alpha}t_{-\alpha} = 1\) 对一切 \(\alpha \in \mathbb{R}\) 成立。

在本节以下部分，我们用 \((X_{\alpha})_{\alpha\in\mathbb{R}}\) 表示满足引理 2 的条件的一个族。用 \(\langle\cdot,\cdot\rangle\) 表示 g 上的一个非退化不变对称双线性型。

每个 \(x \in \mathfrak{g}\) 都可以唯一地写成如下形式

\[
x = h + \sum_ {\alpha \in \mathrm{R}} \mu_ {\alpha} X _ {\alpha} \quad (h \in \mathfrak {h}, \mu_ {\alpha} \in k).
\]

两个这种元素的括号可以用下列公式计算：

\[
[ h, X _ {\alpha} ] = \alpha (h) X _ {\alpha}
\]

\[
[ X _ {\alpha}, X _ {\beta} ] = \left\{ \begin{array}{l l} 0 & \text {if} \alpha + \beta \notin \mathrm{R} \cup \{0 \} \\ - H _ {\alpha} & \text {if} \alpha + \beta = 0 \\ \mathrm{N} _ {\alpha \beta} X _ {\alpha + \beta} & \text {if} \alpha + \beta \in \mathrm{R} \end{array} \right.
\]

其中 \(N_{\alpha\beta}\) 是 k 的非零元素。

引理 3. 对一切 \(\alpha \in \mathbb{R}\) ，

\[
\langle X _ {\alpha}, X _ {- \alpha} \rangle = - \frac {1}{2} \langle H _ {\alpha}, H _ {\alpha} \rangle .
\]

事实上，

\[
\begin{array}{c} 2 \langle X _ {\alpha}, X _ {- \alpha} \rangle = \langle \alpha (H _ {\alpha}) X _ {\alpha}, X _ {- \alpha} \rangle = \langle [ H _ {\alpha}, X _ {\alpha} ], X _ {- \alpha} \rangle \\ = \langle H _ {\alpha}, [ X _ {\alpha}, X _ {- \alpha} ] \rangle = - \langle H _ {\alpha}, H _ {\alpha} \rangle . \end{array}
\]

引理 4. 设 \(\alpha, \beta \in \mathrm{R}\) 使 \(\alpha + \beta \in \mathrm{R}\) 。设 \(p\) （相应地 \(q\) ）是最大的整数 \(j\) ，它使 \(\beta + j\alpha \in \mathrm{R}\) （相应地 \(\beta - j\alpha \in \mathrm{R}\) ）。则

\[
\mathrm{N} _ {\alpha , \beta} \mathrm{N} _ {- \alpha , \alpha + \beta} = - p (q + 1)\tag{3}
\]

\[
\mathrm{N} _ {- \alpha , \alpha + \beta} \langle H _ {\beta}, H _ {\beta} \rangle = - \mathrm{N} _ {- \alpha , - \beta} \langle H _ {\alpha + \beta}, H _ {\alpha + \beta} \rangle\tag{4}
\]

\[
\mathrm{N} _ {\alpha , \beta} \mathrm{N} _ {- \alpha , - \beta} = (q + 1) ^ {2}.\tag{5}
\]

设 \(\rho\) 是 \(\mathfrak{sl}(2,k)\) 在 \(\mathfrak{g}\) 上的、由 \(X_{\alpha}\) 定义的表示。元素 \(e = X_{\beta + p\alpha}\) 是权为 \(p + q\) 的本原元（命题 4 (i)）。置

\[
e _ {n} = \frac {(- 1) ^ {n}}{n !} \rho (X _ {-}) ^ {n} e \quad \text { for } n \geq 0.
\]

由 §1 命题 1，

\[
(\mathrm{ad} X _ {\alpha}) e _ {p} = (q + 1) e _ {p - 1}
\]

\[
(\operatorname{ad} X _ {- \alpha}) (\operatorname{ad} X _ {\alpha}) e _ {p} = - p (q + 1) e _ {p}.
\]

这就证明了 (3)，因为 \(e_p\) 是 \(\mathfrak{g}^{\beta}\) 的非零元素。

由于形式 \(\langle \cdot, \cdot \rangle\) 是不变的，我们有

\[
\langle [ X _ {- \alpha}, X _ {\alpha + \beta} ], X _ {- \beta} \rangle = - \langle X _ {\alpha + \beta}, [ X _ {- \alpha}, X _ {- \beta} ] \rangle
\]

所以

\[
\mathrm{N} _ {- \alpha , \alpha + \beta} \langle X _ {\beta}, X _ {- \beta} \rangle = - \mathrm{N} _ {- \alpha , - \beta} \langle X _ {\alpha + \beta}, X _ {- \alpha - \beta} \rangle
\]

由此，借助引理 3，即证明了 (4)。

把 \(\langle\cdot,\cdot\rangle\) 限制到 h 上是非退化的且在 Weyl 群下不变（命题 5）。借助这一限制把 h 与 \(h^{*}\) 等同起来。若 \(\gamma\in R\) ，则 \(H_{\gamma}\) 等同于 \(2\gamma/\langle\gamma,\gamma\rangle\) （第六章 §1 第 1 节 引理 2）；于是，对一切 \(\gamma,\delta\in R\) ，

\[
\frac {\langle \gamma , \gamma \rangle}{\langle \delta , \delta \rangle} = \frac {\langle H _ {\delta} , H _ {\delta} \rangle}{\langle H _ {\gamma} , H _ {\gamma} \rangle}.\tag{6}
\]

现在，由第六章 §1 第 3 节 命题 10，

\[
\frac {\langle \alpha + \beta , \alpha + \beta \rangle}{\langle \beta , \beta \rangle} = \frac {q + 1}{p}\tag{7}
\]

于是，由 (3), (4), (6), (7)，

\[
\begin{array}{r} \mathrm{N} _ {\alpha , \beta} \mathrm{N} _ {- \alpha , - \beta} = - \mathrm{N} _ {\alpha , \beta} \mathrm{N} _ {- \alpha , \alpha + \beta} \frac {\langle H _ {\beta} , H _ {\beta} \rangle}{\langle H _ {\alpha + \beta} , H _ {\alpha + \beta} \rangle} \\ = - \mathrm{N} _ {\alpha , \beta} \mathrm{N} _ {- \alpha , \alpha + \beta} \frac {q + 1}{p} = (q + 1) ^ {2}. \end{array}
\]

定义 3. \((\mathfrak{g},\mathfrak{h})\) 的一个谢瓦莱系是满足以下条件的族 \((X_{\alpha})_{\alpha \in \mathbb{R}}\) ：

\begin{enumerate}
\def\labelenumi{(\roman{enumi})}
\item
  \(X_{\alpha} \in g^{\alpha}\) 对一切 \(\alpha \in R;\) 成立；
\item
  \([X_{\alpha}, X_{-\alpha}] = -H_{\alpha}\) 对一切 \(\alpha \in \mathrm{R}\) 成立；
\item
  从 \(\mathfrak{g}\) 到 \(\mathfrak{g}\) 的线性映射，若它等于 \(-1\) （在 \(\mathfrak{h}\) 上），并且把 \(X_{\alpha}\) 映到 \(X_{-\alpha}\) （对一切 \(\alpha \in \mathrm{R}\) ），则是 \(\mathfrak{g}\) 的一个自同构。
\end{enumerate}

把这一定义推广到 \((\mathfrak{g},\mathfrak{h})\) 是分裂约化的情形是直接的。

我们将证明（§4 第 4 节 命题 5 的推论），\((\mathfrak{g},\mathfrak{h})\) 的谢瓦莱系是存在的。

命题 7. 设 \((X_{\alpha})_{\alpha \in \mathbb{R}}\) 是 \((\mathfrak{g},\mathfrak{h})\) 的一个谢瓦莱系。我们保留引理 4 的记号。则 \(\mathrm{N}_{-\alpha, - \beta} = \mathrm{N}_{\alpha ,\beta}\) ，且 \(\mathrm{N}_{\alpha ,\beta} = \pm (q + 1)\) 对 \(\alpha ,\beta ,\alpha +\beta \in \mathbb{R}\) 成立。

设 \(\varphi\) 是定义 3 (iii) 中考虑的 \(\mathfrak{g}\) 的自同构。则

\[
\begin{array}{r} \mathrm{N} _ {- \alpha , - \beta} X _ {- \alpha - \beta} = [ X _ {- \alpha}, X _ {- \beta} ] = [ \varphi (X _ {\alpha}), \varphi (X _ {\beta}) ] = \varphi ([ X _ {\alpha}, X _ {\beta} ]) \\ = \varphi (\mathrm{N} _ {\alpha , \beta} X _ {\alpha + \beta}) = \mathrm{N} _ {\alpha , \beta} X _ {- \alpha - \beta} \end{array}
\]

于是 \(N_{-\alpha,-\beta}=N_{\alpha,\beta}\) 。而由 (5) 得 \(N_{\alpha,\beta}=\pm(q+1)\) 。

命题 8. 设 \((X_{\alpha})_{\alpha \in \mathbb{R}}\) 是 \((\mathfrak{g},\mathfrak{h})\) 的一个谢瓦莱系。设 \(\mathrm{M}\) 是一个 \(\mathbf{Z}\) -子模（\(\mathfrak{h}\) 的），它包含诸 \(H_{\alpha}\) 且含于 \(R^{\vee}\) 的权群之中。设 \(\mathfrak{g}_{\mathbf{Z}}\) 是由 M 与诸 \(X_{\alpha}\) 生成的 Z-子模（g 的）。则 \(\mathfrak{g}_{\mathbf{Z}}\) 是 g 的一个 Z-李子代数，并且从 \(g_{Z} \otimes_{Z} k\) 到 g 的典则映射是一个同构。

若 \(\alpha, \beta \in \mathbb{R}\) 使 \(\alpha + \beta \in \mathbb{R}\) ，则 \(\mathrm{N}_{\alpha, \beta} \in \mathbf{Z}\) （命题 7）。另一方面，若 \(\alpha \in \mathbb{R}\) 且 \(h \in \mathbb{M}\) ，则 \(\alpha(h) \in \mathbf{Z}\) （第六章 §1 第 9 节）。这就证明了 \(\mathfrak{g}_{\mathbf{Z}}\) 是一个 \(\mathbf{Z}\) -李子代数（\(\mathfrak{g}\) 的）。另一方面，M 是秩为 dim \(\mathfrak{h}\) 的自由交换群（代数 第七章 §3 定理 1），故 \(\mathfrak{g}_{\mathbf{Z}}\) 是秩为 dim \(\mathfrak{g}\) 的自由交换群；这就蕴含最后一个断言。

\section{分裂半单李代数的子代数}

在本段中，我们用 \((\mathfrak{g},\mathfrak{h})\) 表示一个分裂半单李代数，用 R 表示它的根系。

\subsection{在 ad h 下稳定的子代数}

引理 1. 设 \(\mathrm{V}\) 是 \(\mathfrak{g}\) 的向量子空间，\(\mathrm{R(V)}\) 是 \(\alpha \in \mathrm{R}\) 中使 \(\mathfrak{g}^{\alpha} \subset \mathrm{V}\) 者组成的集合。则 \((\mathrm{V} \cap \mathfrak{h}) + \sum_{\alpha \in \mathrm{R(V)}} \mathfrak{g}^{\alpha}\) 是 \(\mathrm{V}\) 的在 \(\operatorname{ad} \mathfrak{h}\) 下稳定的最大向量子空间。

V 的向量子空间 W 在 ad h 下稳定当且仅当

\[
W = (W \cap \mathfrak {h}) + \sum_ {\alpha \in R} (W \cap \mathfrak {g} ^ {\alpha})
\]

（代数 第七章 §2 第 2 节 定理 1 之推论 1）。于是 V 的在 ad \(\mathfrak{h}\) 下稳定的最大向量子空间是 \((\mathrm{V} \cap \mathfrak{h}) + \sum_{\alpha \in \mathrm{R}} (\mathrm{V} \cap \mathfrak{g}^{\alpha})\) 。但 \(\mathrm{V} \cap \mathfrak{g}^{\alpha} = \mathfrak{g}^{\alpha}\) （对 \(\alpha \in \mathrm{R}(\mathrm{V})\) ），而 \(\mathrm{V} \cap \mathfrak{g}^{\alpha} = 0\) （对 \(\alpha \notin \mathrm{R}(\mathrm{V})\) ），因为 \(\dim \mathfrak{g}^{\alpha} = 1\) 。证毕。

对 R 的任意子集 P，置

\[
\mathfrak {g} ^ {\mathrm{P}} = \sum_ {\alpha \in \mathrm{P}} \mathfrak {g} ^ {\alpha} \quad \mathfrak {h} _ {\mathrm{P}} = \sum_ {\alpha \in \mathrm{P}} \mathfrak {h} _ {\alpha}.
\]

若 \(\mathrm{P}\subset \mathrm{R}\) 且 \(\mathbf{Q}\subset \mathbb{R}\) ，我们显然有

\[
\begin{array}{c} {[ \mathfrak {h}, \mathfrak {g} ^ {\mathrm{P}} ] \subset \mathfrak {g} ^ {\mathrm{P}}} \\ {[ \mathfrak {g} ^ {\mathrm{P}}, \mathfrak {g} ^ {\mathrm{Q}} ] = \mathfrak {g} ^ {(\mathrm{P+Q}) \cap \mathrm{R}} + \mathfrak {h} _ {\mathrm{P} \cap (- \mathrm{Q})}.} \end{array}\tag{1}
\]

\begin{enumerate}
\def\labelenumi{(\arabic{enumi})}
\setcounter{enumi}{1}
\tightlist
\item
\end{enumerate}

回顾（第六章 §1 第 7 节 定义 4），R 的子集 P 称为闭的，如果条件 \(\alpha \in P\) ，\(\beta \in P\) ，\(\alpha + \beta \in R\) 蕴含 \(\alpha + \beta \in P\) ，换言之，如果 \((P + P) \cap R \subset P\) 。

引理 2. 设 \(\mathfrak{h}'\) 是 \(\mathfrak{h}\) 的向量子空间，\(\mathrm{P}\) 是 \(\mathrm{R}\) 的子集。则 \(\mathfrak{h}' + \mathfrak{g}^{\mathrm{P}}\) 是 \(\mathfrak{g}\) 的子代数当且仅当 \(\mathrm{P}\) 是 \(\mathrm{R}\) 的闭子集并且

\[
\mathfrak {h} ^ {\prime} \supset \mathfrak {h} _ {\mathrm{P} \cap (- \mathrm{P})}.
\]

事实上，

\[
[ \mathfrak {h} ^ {\prime} + \mathfrak {g} ^ {\mathrm{P}}, \mathfrak {h} ^ {\prime} + \mathfrak {g} ^ {\mathrm{P}} ] = [ \mathfrak {h} ^ {\prime}, \mathfrak {g} ^ {\mathrm{P}} ] + [ \mathfrak {g} ^ {\mathrm{P}}, \mathfrak {g} ^ {\mathrm{P}} ] = \mathfrak {h} _ {\mathrm{P} \cap (- \mathrm{P})} + [ \mathfrak {h} ^ {\prime}, \mathfrak {g} ^ {\mathrm{P}} ] + \mathfrak {g} ^ {(\mathrm{P} + \mathrm{P}) \cap \mathrm{R}}.
\]

因此 \(\mathfrak{h}' + \mathfrak{g}^{\mathrm{P}}\) 是 \(\mathfrak{g}\) 的子代数当且仅当

\[
\mathfrak {h} _ {\mathrm{P} \cap (- \mathrm{P})} \subset \mathfrak {h} ^ {\prime} \text { and } \mathfrak {g} ^ {(\mathrm{P} + \mathrm{P}) \cap \mathrm{R}} \subset \mathfrak {g} ^ {\mathrm{P}}
\]

这就证明了引理。

命题 1. (i) g 的在 ad h 下稳定的子代数正是形如 \(h' + g^{P}\) 的向量子空间，其中 P 是 R 的闭子集，而 \(h'\) 是 h 的包含 \(\mathfrak{h}_{\mathrm{P}\cap(-\mathrm{P})}\) 的向量子空间。

\begin{enumerate}
\def\labelenumi{(\roman{enumi})}
\setcounter{enumi}{1}
\tightlist
\item
  设 \(\mathfrak{h}'\) ，\(\mathfrak{h}''\) 是 \(\mathfrak{h}\) 的向量子空间，\(\mathrm{P},\mathrm{Q}\) 是 \(\mathrm{R}\) 的闭子集，且 \(\mathfrak{h}'\supset \mathfrak{h}_{\mathrm{P}\cap (-\mathrm{P})}\) ，\(\mathfrak{h}''\subset \mathfrak{h}'\) ，\(\mathrm{Q}\subset \mathrm{P}\) 。则 \(\mathfrak{h}'' + \mathfrak{g}^{\mathrm{Q}}\) 是 \(\mathfrak{h}' + \mathfrak{g}^{\mathrm{P}}\) 的理想当且仅当
\end{enumerate}

\[
(P + Q) \cap R \subset Q \text { and } \mathfrak {h} _ {P \cap (- Q)} \subset \mathfrak {h} ^ {\prime \prime} \subset \bigcap_ {\alpha \in P, \alpha \notin Q} \operatorname{Ker} \alpha .
\]

断言 (i) 由引理 1 与引理 2 立即得出。设 \(\mathfrak{h}',\mathfrak{h}''\) ，P, Q 如 (ii)。则

\[
[ \mathfrak {h} ^ {\prime} + \mathfrak {g} ^ {\mathrm{P}}, \mathfrak {h} ^ {\prime \prime} + \mathfrak {g} ^ {\mathrm{Q}} ] = \mathfrak {h} _ {\mathrm{P} \cap (- \mathrm{Q})} + [ \mathfrak {h} ^ {\prime}, \mathfrak {g} ^ {\mathrm{Q}} ] + [ \mathfrak {h} ^ {\prime \prime}, \mathfrak {g} ^ {\mathrm{P}} ] + \mathfrak {g} ^ {(\mathrm{P} + \mathrm{Q}) \cap \mathrm{R}}.
\]

于是，\(\mathfrak{h}'' + \mathfrak{g}^{\mathrm{Q}}\) 是 \(\mathfrak{h}' + \mathfrak{g}^{\mathrm{P}}\) 的理想当且仅当

\[
\mathfrak {h} _ {\mathrm{P} \cap (- \mathrm{Q})} \subset \mathfrak {h} ^ {\prime \prime}, [ \mathfrak {h} ^ {\prime \prime}, \mathfrak {g} ^ {\mathrm{P}} ] \subset \mathfrak {g} ^ {\mathrm{Q}}, \mathfrak {g} ^ {(\mathrm{P} + \mathrm{Q}) \cap \mathrm{R}} \subset \mathfrak {g} ^ {\mathrm{Q}}.
\]

这就蕴含 (ii)。

命题 2. 设 \(\mathfrak{a}\) 是 \(\mathfrak{g}\) 的在 ad \(\mathfrak{h}\) 下稳定的子代数，并设 \(\mathfrak{h}' \subset \mathfrak{h}\) ，\(\mathrm{P} \subset \mathrm{R}\) 使 \(\mathfrak{a} = \mathfrak{h}' + \mathfrak{g}^{\mathrm{P}}\) 。

\begin{enumerate}
\def\labelenumi{(\roman{enumi})}
\item
  设 \(\mathfrak{k}\) 是 \(x\in \mathfrak{h}'\) 中使 \(\alpha (x) = 0\) 对一切 \(\alpha \in \mathrm{P}\cap (-\mathrm{P})\) 成立者组成的集合。\(\mathfrak{a}\) 的根基是 \(\mathfrak{k} + \mathfrak{g}^{\mathrm{Q}}\) ，其中 \(\mathrm{Q}\) 是 \(\alpha \in \mathrm{P}\) 中使 \(-\alpha \notin \mathrm{P}\) 者组成的集合。此外，\(\mathfrak{g}^{\mathrm{Q}}\) 是 \(\mathfrak{a}\) 的一个幂零理想。
\item
  a 半单当且仅当 P = -P 且 \(h' = h_{P}\) 。
\item
  \(\mathfrak{a}\) 可解当且仅当 \(\mathrm{P} \cap (-\mathrm{P}) = \emptyset\) 。这时 \([\mathfrak{a}, \mathfrak{a}] = \mathfrak{g}^{\mathrm{S}}\) ，其中
\end{enumerate}

\[
\mathrm{S} = ((\mathrm{P} + \mathrm{P}) \cap \mathrm{R}) \cup \{\alpha \in \mathrm{P} | \alpha (\mathfrak {h} ^ {\prime}) \neq 0 \}.
\]

\begin{enumerate}
\def\labelenumi{(\roman{enumi})}
\setcounter{enumi}{3}
\item
  \(\mathfrak{a}\) 在 \(\mathfrak{g}\) 中约化当且仅当 \(\mathrm{P} = -\mathrm{P}\) 。
\item
  \(\mathfrak{a}\) 由幂零元素组成当且仅当 \(\mathfrak{h}' = 0\) 。这时 \(\mathrm{P} \cap (-\mathrm{P}) = \emptyset\) ，且 \(\mathfrak{a}\) 是幂零的。
\end{enumerate}

我们证明 (v)。若 \(\mathfrak{a}\) 由幂零元素组成，则 \(\mathfrak{a}\) 显然是幂零的，且 \(\mathfrak{h}' = 0\) ，因为 \(\mathfrak{h}\) 的元素都是半单的。设 \(\mathfrak{h}' = 0\) 。由命题 1 (i)，\(P \cap (-P) = \varnothing\) 。由第六章 §1 第 7 节 命题 22，存在 R 的一个房 C 使 \(P \subset R_{+}(C)\) 。因此，存在具有下列性质的整数 \(n > 0\) ：若 \(\alpha_{1}, \ldots, \alpha_{n} \in P\) 且 \(\beta \in R \cup \{0\}\) ，则

\[
\alpha_ {1} + \dots + \alpha_ {n} + \beta \notin \mathrm{R} \cup \{0 \}.
\]

这蕴含 \(\mathfrak{g}^{\mathrm{P}}\) 的每个元素都是幂零的，于是得到 (v)。

我们证明 (iii)。若 \(\mathrm{P} \cap (-\mathrm{P}) = \varnothing\) ，则 \(\mathfrak{g}^{\mathrm{P}}\) 是 \(\mathfrak{g}\) 的子代数（命题 1 (i)），且由 (v) 是幂零的。现在

\[
[ \mathfrak {a}, \mathfrak {a} ] = [ \mathfrak {h} ^ {\prime}, \mathfrak {g} ^ {\mathrm{P}} ] + [ \mathfrak {g} ^ {\mathrm{P}}, \mathfrak {g} ^ {\mathrm{P}} ] = [ \mathfrak {h} ^ {\prime}, \mathfrak {g} ^ {\mathrm{P}} ] + \mathfrak {g} ^ {(\mathrm{P} + \mathrm{P}) \cap \mathrm{R}} \subset \mathfrak {g} ^ {\mathrm{P}},
\]

于是 \(\mathfrak{a}\) 可解，且 \([\mathfrak{a},\mathfrak{a}]\) 由命题中的公式给出。若 \(\mathrm{P}\cap (-\mathrm{P})\neq \emptyset\) ，取 \(\alpha \in \mathrm{P}\) 使 \(-\alpha \in \mathrm{P}\) 。则 \(\mathfrak{h}_{\alpha} + \mathfrak{g}^{\alpha} + \mathfrak{g}^{-\alpha}\) 是 \(\mathfrak{a}\) 的一个单子代数，故 \(\mathfrak{a}\) 不可解。

我们证明 (i)。由于 P 是闭的，\((\mathrm{P} + \mathrm{Q}) \cap \mathrm{R} \subset \mathrm{P}\) 。若 \(\alpha \in P, \beta \in Q\) 且 \(\alpha + \beta \in R\) ，则不可能有 \(\alpha + \beta \in -P\) ，因为，P 既为闭的，这将蕴含 \(-\beta = -(\alpha + \beta) + \alpha \in P\) ，而 \(\beta \in Q\) ；于是，\((\mathrm{P} + \mathrm{Q}) \cap \mathrm{R} \subset \mathrm{Q}\) 。这证明了 \(g^{Q}\) 是 a 的理想，且由 (v) 是幂零的。我们有 \(\mathrm{P} \cap (-\mathrm{Q}) = \varnothing\) ，且 \(\mathrm{P} \cap (-\mathrm{P}) = \mathrm{P} \cap \complement Q\) ，故 \(\mathfrak{h}_{\mathrm{P} \cap (-\mathrm{Q})} \subset \mathfrak{k} \subset \bigcap_{\alpha \in \mathrm{P}, \alpha \notin Q} \operatorname{Ker} \alpha\) 。由命题 1 (ii)，

\(\mathfrak{k} + \mathfrak{g}^{\mathrm{Q}}\) 是 \(\mathfrak{a}\) 的理想。由于 \(\mathrm{Q} \cap (-\mathrm{Q}) = \emptyset\) ，由 (iii) 这个理想是可解的。因此它含于根基 \(\mathfrak{r}\) （\(\mathfrak{a}\) 的）。由于 \(\mathfrak{r}\) 在 \(\mathfrak{a}\) 的每个导子下稳定，\(\mathfrak{r}\) 在 ad \(\mathfrak{h}\) 下稳定。因此存在子集 \(S\) （\(P\) 的），使 \(\mathfrak{r} = (\mathfrak{r} \cap \mathfrak{h}) + \mathfrak{g}^{\mathrm{S}}\) 。设 \(\alpha \in S\) 且 \(-\alpha \in P\) 。则 \(\mathfrak{h}_{\alpha} = [\mathfrak{g}^{\alpha}, \mathfrak{g}^{-\alpha}] \subset \mathfrak{r}\) ，故 \(\mathfrak{g}^{-\alpha} = [\mathfrak{h}_{\alpha}, \mathfrak{g}^{-\alpha}] \subset \mathfrak{r} = 0\) ，从而 \(-\alpha \in S\) ；由 (iii)，这与 \(\mathfrak{r}\) 可解矛盾。因而 \(S \subset Q\) 。最后，若 \(x \in \mathfrak{r} \cap \mathfrak{h}\) 且 \(\alpha \in P \cap (-P)\) ，则 \([x, \mathfrak{g}^{\alpha}] \subset \mathfrak{g}^{\alpha} \cap \mathfrak{r} = 0\) ，故 \(\alpha(x) = 0\) ；这表明 \(x \in \mathfrak{k}\) 。于是 \(\mathfrak{r} \subset \mathfrak{k} + \mathfrak{g}^{\mathrm{Q}}\) ，(i) 的证明完毕。

我们证明 (iv)。由 (i)，\(\mathfrak{a}\) 在 \(\mathfrak{g}\) 上的伴随表示是半单的，当且仅当 \(\mathrm{ad}_{\mathfrak{g}}x\) 对一切 \(x\in \mathfrak{k} + \mathfrak{g}^{\mathrm{Q}}\) 是半单的（第一章 §6 第 5 节 定理 4）；由 (v)，此事成立当且仅当 \(\mathrm{Q} = \varnothing\) ，换言之 \(\mathrm{P} = -\mathrm{P}\) 。

我们证明 (ii)。若 \(\mathfrak{a}\) 半单，则由 (i) 有 \(P = -P\) ，故 \(\mathfrak{h}_{\mathrm{P}} \subset \mathfrak{h}'\) ；再者，\(\mathfrak{a} = [\mathfrak{a}, \mathfrak{a}] \subset \mathfrak{h}_{\mathrm{P}} + \mathfrak{g}^{\mathrm{P}}\) ，从而 \(\mathfrak{h}' = \mathfrak{h}_{\mathrm{P}}\) 。若 \(P = -P\) 且 \(\mathfrak{h}' = \mathfrak{h}_{\mathrm{P}}\) ，则 \(\mathfrak{a}\) 由 (iv) 是约化的，且 \(\mathfrak{a} = \sum_{\alpha \in P} \mathfrak{s}_{\alpha}\) ，故 \(\mathfrak{a} = [\mathfrak{a}, \mathfrak{a}]\) ，\(\mathfrak{a}\) 是半单的。

命题 3. 设 \(\mathfrak{a}\) 是 \(\mathfrak{g}\) 的在 \(\operatorname{ad}(\mathfrak{h})\) 下稳定的半单子代数，并设 \(\mathrm{P}\) 是 \(\mathrm{R}\) 的子集，它使 \(\mathfrak{a} = \mathfrak{h}_{\mathrm{P}} + \mathfrak{g}^{\mathrm{P}}\) 。

\begin{enumerate}
\def\labelenumi{(\roman{enumi})}
\item
  \(\mathfrak{h}_{\mathrm{P}}\) 是 \(\mathfrak{a}\) 的一个分裂嘉当子代数。
\item
  \((\mathfrak{a},\mathfrak{h}_{\mathrm{P}})\) 的根系是诸元素到 \(\mathfrak{h}_{\mathrm{P}}\) 上的限制所成的集合，这些元素取自 \(\mathrm{P}\) 。
\end{enumerate}

由于 \(h_{P}\) 在 ad h 下稳定，它在 a 中的正规化子在 ad h 下稳定，因而具有形式 \(h_{P} + g^{Q}\) ，其中 \(Q \subset P\) （引理 1）。若 \(\alpha \in Q\) ，

\[
\mathfrak {g} ^ {\alpha} = \left[ \mathfrak {h} _ {\alpha}, \mathfrak {g} ^ {\alpha} \right] \subset \left[ \mathfrak {h} _ {\mathrm{P}}, \mathfrak {g} ^ {\alpha} \right] \subset \mathfrak {h} _ {\mathrm{P}},
\]

这是荒谬的。于是 \(Q = \varnothing\) ，\(h_{P}\) 是它在 a 中自己的正规化子。这证明了 \(h_{P}\) 是 a 的一个嘉当子代数。若 \(x \in h_{P}\) ，则 \(ad_{g}x\) ，从而更有 \(ad_{a}x\) ，是可三角化的。于是 (i) 得证，而 (ii) 是明显的。

由命题 1 (i)，g 的包含 h 的子代数正是那些集合 \(h + g^{P}\) ，其中 P 是 R 的闭子集。由第七章 §3 命题 3，\(h + g^{P}\) 的每个嘉当子代数都是 g 的嘉当子代数。

命题 4. 设 \(\mathfrak{a}\) 是 \(\mathfrak{g}\) 的包含 \(\mathfrak{h}\) 的子代数，\(x\) 是 \(\mathfrak{a}\) 的元素，\(s\) 与 \(n\) 是它的半单分量与幂零分量。则 \(s \in \mathfrak{a}\) 且 \(n \in \mathfrak{a}\) 。

我们有 \((\operatorname{ad} x)\mathfrak{a} \subset \mathfrak{a}\) ，故 \((\operatorname{ad} s)\mathfrak{a} \subset \mathfrak{a}\) 且 \((\operatorname{ad} n)\mathfrak{a} \subset \mathfrak{a}\) 。由于 a 是它在 g 中自己的正规化子（第七章 §2 第 1 节 命题 4 之推论 4），\(s \in a\) 且 \(n \in a\) 。

命题 5. 设 P 是 R 的闭子集。

\begin{enumerate}
\def\labelenumi{(\roman{enumi})}
\item
  \(\mathfrak{h} + \mathfrak{g}^{\mathrm{P}}\) 可解当且仅当 \(\mathrm{P} \cap (-\mathrm{P}) = \varnothing\) 。这时，\([\mathfrak{h} + \mathfrak{g}^{\mathrm{P}}, \mathfrak{h} + \mathfrak{g}^{\mathrm{P}}] = \mathfrak{g}^{\mathrm{P}}\) 。
\item
  \(\mathfrak{h} + \mathfrak{g}^{\mathrm{P}}\) 约化当且仅当 \(\mathrm{P} = -\mathrm{P}\) 。
\end{enumerate}

断言 (i) 由命题 2 (iii) 得出。若 \(\mathrm{P} = -\mathrm{P}\) ，则 \(\mathfrak{h} + \mathfrak{g}^{\mathrm{P}}\) 是约化的（命题 2 (iv)）。设 \(\mathfrak{a} = \mathfrak{h} + \mathfrak{g}^{\mathrm{P}}\) 是约化的。则

\[
\mathfrak {g} ^ {\mathrm{P}} = \left[ \mathfrak {h}, \mathfrak {g} ^ {\mathrm{P}} \right] \subset \left[ \mathfrak {a}, \mathfrak {a} \right] \subset \mathfrak {h} + \mathfrak {g} ^ {\mathrm{P}},
\]

于是 \([a, a]\) 具有形式 \(h' + g^{P}\) ，其中 \(h' \subset h\) ；由于 \([a, a]\) 半单，故 P = -P （命题 2 (ii)）。

\subsection{理想}

命题 6. 设 \(\mathrm{R}_1, \ldots, \mathrm{R}_p\) 是 \(\mathbb{R}\) 的不可约分支。对 \(i = 1, \ldots, p\) ，置 \(\mathfrak{g}_i = \mathfrak{h}_{\mathrm{R}_i} + \mathfrak{g}^{\mathrm{R}_i}\) 。则 \(\mathfrak{g}_1, \ldots, \mathfrak{g}_p\) 是 \(\mathfrak{g}\) 的单分量。

诸 \(g_{i}\) 是 g 的理想（命题 1 (ii)）。显然 g 是诸 \(g_{i}\) 的直和，因而是诸 \(g_{i}\) 的积。设 a 与 b 是 g 的互补理想。则 a 与 b 都是半单的且在 ad h 下稳定，故存在 R 的子集 P, Q 使 \(a = h_{P} + g^{P}\) ，\(b = h_{Q} + g^{Q}\) 。于是 \(h_{P}\) ，\(h_{Q}\) 关于 Killing 型在 h 中互为正交补，故 P 与 Q 是 R 的不可约分支的并。这证明了诸 \(g_{i}\) 是 g 的极小理想。

推论 1. g 单当且仅当 R 不可约（换言之，它的邓金图是连通的）。

这由命题 6 得出。

李代数 \(\mathfrak{a}\) 称为绝对单的，如果对每个扩张 \(k'\) （\(k\) 的），\(k'\) -李代数 \(\mathfrak{a}_{(k')}\) 都是单的。

推论 2. 可分裂单李代数是绝对单的。

这由推论 1 得出。

若 g 是 \(A_{l}\) 型（ \(l \geq 1\) ）、\(B_{l}\) 型（ \(l \geq 1\) ）、\(C_{l}\) 型（ \(l \geq 1\) ）或 \(D_{l}\) 型（ \(l \geq 3\) ）的，则称 g 为经典可分裂单李代数。若 g 是 \(E_{6}\) ，\(E_{7}\) ，\(E_{8}\) ，\(F_{4}\) 或 \(G_{2}\) 型的，则称 g 为例外型可分裂单李代数。

\subsection{Borel 子代数}

命题 7. 设 \(\mathfrak{b} = \mathfrak{h} + \mathfrak{g}^{\mathrm{P}}\) 是 \(\mathfrak{g}\) 的包含 \(\mathfrak{h}\) 的子代数。下列条件等价：

\begin{enumerate}
\def\labelenumi{(\roman{enumi})}
\item
  \(\mathfrak{b}\) 是 \(\mathfrak{g}\) 的极大可解子代数；
\item
  存在 R 的房 C 使 \(P = R_{+}(C)\) ；
\item
  \(\mathrm{P} \cap (-\mathrm{P}) = \varnothing\) 且 \(\mathrm{P} \cup (-\mathrm{P}) = \mathrm{R}\) 。
\item
  \(\Longrightarrow\) (ii)：若 \(\mathfrak{b}\) 可解，则 \(P \cap (-P) = \varnothing\) 。于是存在房 \(C\) （\(R\) 中的）使 \(P \subset R_{+}(C)\) （第六章 §1 第 7 节 命题 22）。这时 \(\mathfrak{h} + \mathfrak{g}^{R_{+}(C)}\) 是 \(\mathfrak{g}\) 的包含 \(\mathfrak{b}\) 的可解子代数，故若 \(\mathfrak{b}\) 极大则等于 \(\mathfrak{b}\) 。
\item
  \(\Longrightarrow\) (iii)：这是明显的。
\item
  \(\Longrightarrow\) (i)：设 \(P \cap (-P) = \varnothing\) 且 \(P \cup (-P) = R\) 。则 b 可解。设 \(b'\) 是 g 的包含 b 的可解子代数。存在 R 的子集 Q 使 \(b' = h + g^{Q}\) 。则 \(Q \cap (-Q) = \varnothing\) 且 \(Q \supset P\) ，故 Q = P 且 \(b' = b\) 。
\end{enumerate}

定义 1. g 的包含 h 且满足命题 7 中等价条件的子代数称为 (g, h) 的 Borel 子代数。

子代数 \(\mathfrak{b}\) （可分裂代数 \(\mathfrak{g}\) 的）称为 \(\mathfrak{g}\) 的 Borel 子代数，如果存在分裂嘉当子代数 \(\mathfrak{h}'\) （\(\mathfrak{g}\) 的）使 \(\mathfrak{b}\) 是 \((\mathfrak{g},\mathfrak{h}')\) 的 Borel 子代数。

设 \((\mathfrak{g},\mathfrak{h})\) 是分裂约化李代数。设 \(\mathfrak{g}=\mathfrak{c}\times\mathfrak{s}\) ，其中 c 交换而 s 半单。g 的形如 \(\mathfrak{c}\times\mathfrak{b}\) 的子代数（其中 b 是 \((\mathfrak{s},\mathfrak{h}\cap\mathfrak{s})\) 的 Borel 子代数）称为 \((\mathfrak{g},\mathfrak{h})\) 的 Borel 子代数。

采用命题 7 的记号，我们也说 b 是由 h 与 C（或由 h 与对应于 C 的 R 的基）所定义的 g 的 Borel 子代数。

注. 把 \(\mathrm{R}_{+}(\mathrm{C})\) 对应于 R 的每个房 C 的映射是单射（第六章 §1 第 7 节 命题 20 之推论 1）。因此，\(\mathrm{C} \mapsto \mathfrak{h} + \mathfrak{g}^{\mathrm{R}_{+}(\mathrm{C})}\) 是从 R 的房之集到 \((\mathfrak{g}, \mathfrak{h})\) 的 Borel 子代数之集上的双射。于是，\((\mathfrak{g}, \mathfrak{h})\) 的 Borel 子代数的个数等于 R 的 Weyl 群的阶（第六章 §1 第 5 节 定理 2）。

命题 8. 设 \(\mathfrak{b}\) 是 \(\mathfrak{g}\) 的子代数，\(k'\) 是 \(k\) 的扩张。则 \(\mathfrak{b} \otimes_k k'\) 是 \((\mathfrak{g} \otimes_k k', \mathfrak{h} \otimes_k k')\) 的 Borel 子代数当且仅当 \(\mathfrak{b}\) 是 \((\mathfrak{g}, \mathfrak{h})\) 的 Borel 子代数。

这由命题 7 的条件 (iii) 立即得出。

命题 9. 设 \(\mathfrak{b}\) 是由 R 的房 C 所定义的 \((\mathfrak{g},\mathfrak{h})\) 的 Borel 子代数。设 \(\mathfrak{n} = \mathfrak{g}^{\mathrm{R}_{+}(\mathrm{C})} = \sum_{\alpha \in \mathrm{R},\alpha >0}\mathfrak{g}^{\alpha}\) 。设 \(l = \dim \mathfrak{h}\) 。

\begin{enumerate}
\def\labelenumi{(\roman{enumi})}
\item
  若 \(h \in \mathfrak{h}\) 且 \(x \in \mathfrak{n}\) ，则 \(\mathrm{ad}_{\mathfrak{g}}(h + x)\) 的特征多项式是 \(\mathrm{T}^l \prod_{\alpha \in \mathbb{R}} (\mathrm{T} - \alpha(h))\) 。
\item
  \(\mathfrak{b}\) 的最大幂零理想等于 \(\mathfrak{n}\) ，也等于 \([\mathfrak{b},\mathfrak{b}]\) 。它也是 \(\mathfrak{b}\) 中在 \(\mathfrak{g}\) 内幂零的元素组成的集合。
\item
  设 \(\mathrm{B}\) 是 \(\mathbb{R}\) 的、对应于 \(\mathbb{C}\) 的基。对一切 \(\alpha \in \mathbb{B}\) ，设 \(X_{\alpha}\) 是 \(\mathfrak{g}^{\alpha}\) 的非零元素。则 \((X_{\alpha})_{\alpha \in \mathbb{B}}\) 生成李代数 \(\mathfrak{n}\) 。我们有 \([\mathfrak{n}, \mathfrak{n}] = \sum_{\alpha \in \mathbb{R}, \alpha > 0, \alpha \notin \mathbb{B}} \mathfrak{g}^{\alpha}\) 。
\end{enumerate}

在 \(\mathfrak{h}_{\mathbf{Q}}^{*}\) 上存在一个全序，它与它的向量空间结构相容，且使 \(\mathrm{R}_{+}(\mathrm{C})\) 的元素都 \(>0\) （第六章 §1 第 7 节）。设 \(h, x\) 如 (i)，并设 \(y \in \mathfrak{g}^{\alpha}\) 。则 \([h + x, y] = \alpha(h)y + z\) ，其中 \(z \in \sum_{\beta > \alpha} \mathfrak{g}^{\beta}\) 。于是，关于 \(\mathfrak{g}\) 的一个适当的基，\(\operatorname{ad}_{\mathfrak{g}}(h + x)\) 的矩阵具有下列性质：

\begin{enumerate}
\def\labelenumi{\arabic{enumi})}
\item
  它是下三角的；
\item
  矩阵的对角元是数 0（l 次）以及诸 \(\alpha(h)\) ，其中 \(\alpha \in \mathbb{R}\) 。
\end{enumerate}

这就证明了 (i)。它还表明 \(\operatorname{ad}_{\mathfrak{b}}(h+x)\) 的特征多项式是 \(\mathrm{T}^{l}\prod_{\alpha\in\mathrm{R}_{+}(\mathrm{C})}(\mathrm{T}-\alpha(h))\) 。由前面的论证，b 中在 g 内幂零的元素之集以及 b 的最大幂零理想都等于 n。~由命题 5 (i) 有 \(n=[b,b]\) 。最后，断言 (iii) 由 §2 命题 4 (ii) 以及第六章 §1 第 6 节 命题 19 得出。

推论. 设 \(\mathfrak{b}\) 是 \(\mathfrak{g}\) 的 Borel 子代数。

\begin{enumerate}
\def\labelenumi{(\roman{enumi})}
\item
  b 的每个嘉当子代数都是 g 的分裂嘉当子代数。
\item
  若 \(\mathfrak{h}_1, \mathfrak{h}_2\) 是 \(\mathfrak{b}\) 的嘉当子代数，则存在 \(x \in [\mathfrak{b}, \mathfrak{b}]\) 使 \(e^{\mathrm{ad}_{\mathfrak{g}} x} \mathfrak{h}_1 = \mathfrak{h}_2\) 。
\end{enumerate}

断言 (i) 由命题 9 (i) 与第七章 §3 第 3 节 命题 3 得出。断言 (ii) 由命题 9 (ii) 与第七章 §3 第 4 节 定理 3 得出。

命题 10. 设 \(\mathfrak{b},\mathfrak{b}'\) 是 \(\mathfrak{g}\) 的 Borel 子代数。则存在 \(\mathfrak{g}\) 的分裂嘉当子代数含于 \(\mathfrak{b} \cap \mathfrak{b}'\) 。

设 \(\mathfrak{h}\) 是 \(\mathfrak{b}\) 的嘉当子代数，\(\mathfrak{n} = [\mathfrak{b},\mathfrak{b}],\mathfrak{n}' = [\mathfrak{b}',\mathfrak{b}']\) ，\(\mathfrak{p} = \mathfrak{b}\cap \mathfrak{b}'\) ，并设 \(\mathfrak{s}\) 是 \(\mathfrak{g}\) 的、与 \(\mathfrak{b} + \mathfrak{b}'\) 互补的一个向量子空间。用 \(\mathfrak{s}^{\perp},\mathfrak{b}^{\perp},\mathfrak{b}'^{\perp}\) 表示 \(\mathfrak{s},\mathfrak{b},\mathfrak{b}^{\prime}\) 关于 \(\mathfrak{g}\) 的 Killing 型的正交补。置 \(l = \dim \mathfrak{h}, n = \dim \mathfrak{n}, p = \dim \mathfrak{p}\) 。则 \(\dim \mathfrak{b} = \dim \mathfrak{b}^{\prime} = l + n\) ，

\[
\dim \mathfrak {s} ^ {\perp} = \dim (\mathfrak {b} + \mathfrak {b} ^ {\prime}) = 2 (l + n) - p,
\]

于是

\[
\begin{array}{c} \dim (\mathfrak {s} ^ {\perp} \cap \mathfrak {p}) \geq \dim \mathfrak {s} ^ {\perp} + \dim \mathfrak {p} - \dim \mathfrak {g} \\ = 2 (l + n) - p + p - (l + 2 n) = l. \end{array}\tag{3}
\]

由 §2 第 2 节 命题 1，\(\mathfrak{n} \subset \mathfrak{b}^{\perp}\) ，\(\mathfrak{n}' \subset \mathfrak{b}'^{\perp}\) 。\(\mathfrak{p} \cap \mathfrak{n}\) 的元素在 \(\mathfrak{g}\) 内是幂零的（命题 9 (ii)），并且属于 \(\mathfrak{b}'\) ，从而属于 \(\mathfrak{n}'\) （命题 9 (ii)）。因此，\(\mathfrak{p} \cap \mathfrak{n} \subset \mathfrak{n} \cap \mathfrak{n}' \subset \mathfrak{b}^{\perp} \cap \mathfrak{b}'^{\perp}\) ，故 \(\mathfrak{s}^{\perp} \cap \mathfrak{p} \cap \mathfrak{n} = 0\) 。鉴于 (3)，我们看到 \(\mathfrak{s}^{\perp} \cap \mathfrak{p}\) 是 \(\mathfrak{n}\) 在 \(\mathfrak{b}\) 中的补。设 \(z\) 是 \(\mathfrak{h}\) 中在 \(\mathfrak{g}\) 内正则的一个元素；存在 \(y \in \mathfrak{n}\) 使 \(y + z \in \mathfrak{s}^{\perp} \cap \mathfrak{p}\) ；由命题 9 (i)，\(\operatorname{ad}_{\mathfrak{g}}(y + z)\) 与 \(\operatorname{ad}_{\mathfrak{g}}z\) 有相同的特征多项式，故 \(x = y + z\) 在 \(\mathfrak{g}\) 中正则，从而更有在 \(\mathfrak{b}\) 中与 \(\mathfrak{b}'\) 中正则（第七章 §2 第 2 节 命题 9）。由于 \(\mathfrak{g}, \mathfrak{b}, \mathfrak{b}'\) 有相同的秩，\(\mathfrak{b}^{0}(x) = \mathfrak{g}^{0}(x) = \mathfrak{b}'^{0}(x)\) 同时是 \(\mathfrak{b}\) 的、\(\mathfrak{g}\) 的与 \(\mathfrak{b}'\) 的嘉当子代数（第七章 §3 第 3 节 定理 2）。最后，\(\mathfrak{g}\) 的这个嘉当子代数由命题 9 的推论是分裂的。

推论. 群 \(\operatorname{Aut}_{e}(\mathfrak{g})\) 传递地作用于偶 \((\mathfrak{t},\mathfrak{b})\) 组成的集合上，其中 t 是 g 的分裂嘉当子代数，b 是 \((\mathfrak{g},\mathfrak{t})\) 的 Borel 子代数。

设 \((\mathfrak{t}_1, \mathfrak{b}_1)\) 与 \((\mathfrak{t}_2, \mathfrak{b}_2)\) 是两个这样的偶。存在分裂嘉当子代数 \(\mathfrak{t}\) （\(\mathfrak{g}\) 的）含于 \(\mathfrak{b}_1 \cap \mathfrak{b}_2\) （命题 10）。由命题 9 的推论，我们归结到 \(\mathfrak{t}_1 = \mathfrak{t}_2 = \mathfrak{t}\) 的情形。设 \(S\) 是 \((\mathfrak{g}, \mathfrak{t})\) 的根系。存在基 \(B_1, B_2\) （\(S\) 的）使 \(\mathfrak{b}_i\) 对应于 \(B_i\) （ \(i = 1, 2\) ），且存在 \(s \in W(S)\) 把 \(B_1\) 变为 \(B_2\) 。最后，存在 \(a \in \operatorname{Aut}_e(\mathfrak{g})\) 使 \(a|t = s\) （ \(\S 2\) 第 2 节 定理 2 的推论）。则 \(a(t) = t\) 且 \(a(\mathfrak{b}_1) = \mathfrak{b}_2\) 。

\subsection{抛物子代数}

命题 11. 设 \(\mathfrak{p} = \mathfrak{h} + \mathfrak{g}^{\mathrm{P}}\) 是 \(\mathfrak{g}\) 的包含 \(\mathfrak{h}\) 的子代数。下列条件等价：

\begin{enumerate}
\def\labelenumi{(\roman{enumi})}
\item
  p 包含 \((\mathfrak{g}, \mathfrak{h})\) 的一个 Borel 子代数；
\item
  存在 R 的房 C 使 \(\mathrm{P} \supset \mathrm{R}_{+}(\mathrm{C})\) ；
\item
  P 是抛物的，换言之（第六章 §1 第 7 节 定义 4），\(\mathrm{P} \cup (-\mathrm{P}) = \mathrm{R}\) 。
\end{enumerate}

条件 (i) 与 (ii) 的等价性由命题 7 得出。条件 (ii) 与 (iii) 的等价性由第六章 §1 第 7 节 命题 20 得出。

定义 2. g 的包含 h 且满足命题 11 的等价条件的子代数称为 \((\mathfrak{g}, \mathfrak{h})\) 的抛物子代数。g 的抛物子代数是指 \((\mathfrak{g}, \mathfrak{h}')\) 的抛物子代数，其中 \(h'\) 是 g 的一个分裂嘉当子代数。

这一定义立即推广到 \((\mathfrak{g}, \mathfrak{h})\) 是分裂约化李代数的情形。

注. 设 B 是 R 的一个基，b 是对应的 Borel 子代数。若 \(\Sigma \subset B\) ，用 \(Q_{\Sigma}\) 表示 \(\Sigma\) 中元素的具系数 \(\leq 0\) 的线性组合所成之根组成的集合；置 \(\mathfrak{p}(\Sigma) = \mathrm{R}_{+}(\mathrm{B}) \cup \mathrm{Q}_{\Sigma}\) 与 \(\mathfrak{p}_{\Sigma} = \mathfrak{h} \oplus \mathfrak{g}^{\mathrm{P}(\Sigma)}\) 。由第六章 §1 第 7 节 引理 3 与命题 20，\(p_{\Sigma}\) 是包含 b 的抛物子代数，且 g 的每个包含 b 的抛物子代数都是这样得到的。

引理 3. 设 V 是有限维实向量空间，S 是 \(V^{*}\) 中的一个根系，P 是 S 的抛物子集组成的集合；设 H 是诸 Ker \(\alpha\) （\(\alpha \in S\) ）组成的集合，F 是 V 关于 H 的面组成的集合（第五章 §1 第 2 节 定义 1）。

若 \(\mathrm{P} \in \mathcal{P}\) ，设 \(\overline{\mathrm{F}}(\mathrm{P})\) 是 \(v \in \mathrm{V}\) 中使 \(\alpha(v) \geq 0\) 对一切 \(\alpha \in \mathrm{P}\) 成立者组成的集合。若 \(\mathrm{F} \in \mathcal{F}\) ，设 \(\mathrm{P}(\mathrm{F})\) 是 \(\alpha \in \mathrm{R}\) 中使 \(\alpha(v) \geq 0\) 对一切 \(v \in \mathrm{F}\) 成立者组成的集合。

则 \(\mathrm{F} \mapsto \mathrm{P}(\mathrm{F})\) 是从 \(\mathcal{F}\) 到 \(\mathcal{P}\) 的一个双射；对一切 \(\mathrm{F} \in \mathcal{F}\) ，\(\overline{\mathrm{F}}(\mathrm{P}(\mathrm{F}))\) 是 \(\mathrm{F}\) 的闭包。

\begin{enumerate}
\def\labelenumi{\alph{enumi})}
\tightlist
\item
  设 \(P \in P\) 。存在 S 的一个房 C 以及一个子集 \(\Sigma\) （基 \(\mathrm{B}(\mathrm{C})\) 的），使 \(\mathrm{P} = \mathrm{S}_{+}(\mathrm{C}) \cup \mathrm{Q}\) ，其中 Q 是 \(\Sigma\) 中元素的具非正整系数的线性组合组成的集合（第六章 §1 第 7 节 命题 20）。置
\end{enumerate}

\[
\mathrm{B} (\mathrm{C}) = \{\alpha_ {1}, \dots , \alpha_ {l} \}, \Sigma = \{\alpha_ {1}, \dots , \alpha_ {m} \}.
\]

若 \(v \in \mathrm{V}\) ，我们有下列等价关系：

\[
\begin{array}{c} \alpha (v) \geq 0 \text {   for   all   } \alpha \in \mathrm{P} \\ \Longleftrightarrow \alpha_ {1} (v) \geq 0, \ldots \alpha_ {l} (v) \geq 0, \alpha_ {1} (v) \leq 0, \ldots , \alpha_ {m} (v) \leq 0 \\ \Longleftrightarrow \alpha_ {1} (v) = \dots = \alpha_ {m} (v) = 0, \alpha_ {m + 1} (v) \geq 0, \ldots , \alpha_ {l} (v) \geq 0, \end{array}
\]

故 \(\overline{\mathrm{F}}(\mathrm{P})\) 是集合

\[
\{v \in \mathrm{V} \mid \alpha_ {1} (v) = \dots = \alpha_ {m} (v) = 0, \alpha_ {m + 1} (v) > 0, \dots , \alpha_ {l} (v) > 0 \},
\]

的闭包，此集合是相对于 H 的一个面 F ，因为 S 的每个元素都是 \(\alpha_{1},\ldots,\alpha_{l}\) 的线性组合，其系数或者全 \(\geq0\) 或者全 \(\leq0\) 。此外，若 \(\beta=u_{1}\alpha_{1}+\cdots+u_{l}\alpha_{l}\in S\) ，

\[
\begin{array}{c} \beta \in \mathrm{P} (\mathrm{F}) \Longleftrightarrow u _ {m + 1} \geq 0, \ldots , u _ {l} \geq 0 \\ \Longleftrightarrow \beta \in \mathrm{S} _ {+} (\mathrm{C}) \text {   or   } (- \beta \in \mathrm{S} _ {+} (\mathrm{C}) \text {   and   } u _ {m + 1} = \ldots = u _ {l} = 0) \\ \Longleftrightarrow \beta \in \mathrm{S} _ {+} (\mathrm{C}) \cup \mathrm{Q} = \mathrm{P}, \end{array}
\]

故 P(F) = P。

\begin{enumerate}
\def\labelenumi{\alph{enumi})}
\setcounter{enumi}{1}
\tightlist
\item
  设 \(F \in F\) 。显然 \(P(F) \in \mathcal{P}\) 。另一方面，F 含于相对于 \(\mathcal{H}(\text{Chap. V}, \S1, \text{no. 3, formulas (6)}),\) 的一个房的闭包之中，因而是相对于该房的壁所成集合的一个面（第五章 \(\S1\) 第 4 节 命题 9）。于是，\(\overline{F}\) 具有形状 \(\{v \in V \mid \alpha(v) \geq 0 \text{ 对一切 } \alpha \in T\}\) ，其中 T 是 S 的一个子集，而显然可取它等于 \(P(F)\) 。这样，\(\overline{F} = \overline{F}(P(F))\) 。证毕。
\end{enumerate}

若 \(P \in P\) ，使 \(P = P(F)\) 的面 F 称为与 P 相伴的面；我们把它记作 \(F(P)\) 。我们把这些约定推广到 \((\mathfrak{g}, \mathfrak{h})\) 是分裂约化的情形。

命题 12. 设 \(\mathcal{H}\) 是 \(\mathfrak{h}_{\mathbf{R}}\) 的超平面组成的集合，这些超平面由 R 中根的核构成。设 \(\mathcal{F}\) 是 \(\mathfrak{h}_{\mathbf{R}}\) 相对于 \(\mathcal{H}\) 的面组成的集合。设 \(\mathcal{S}\) 是 \((\mathfrak{g},\mathfrak{h})\) 的抛物子代数组成的集合。对每个 \(\mathfrak{p} = \mathfrak{h} + \mathfrak{g}^{\mathrm{P}}\in \mathcal{S}\) ，设 \(\mathrm{F}(\mathfrak{p})\) 是与 P 相伴的面。则 \(\mathfrak{p}\mapsto \mathrm{F}(\mathfrak{p})\) 是从 \(\mathcal{S}\) 到 \(\mathcal{F}\) 的一个双射。若 \(\mathfrak{p}_1,\mathfrak{p}_2\in \mathcal{P}\) ，

\[
\mathfrak {p} _ {1} \supset \mathfrak {p} _ {2} \Longleftrightarrow \mathrm{F} (\mathfrak {p} _ {1}) \subset \overline {{\mathrm{F} (\mathfrak {p} _ {2})}}.
\]

这由引理 3 立即得出。

例. 对应于 \((\mathfrak{g}, \mathfrak{h})\) 的包含一个 Borel 代数 b 的抛物子代数的那些面，就是含于与 b 相伴的房的闭包之中的那些面（参见上面的注）。

命题 13. 设 \(\mathfrak{p} = \mathfrak{h} + \mathfrak{g}^{\mathrm{P}}\) 是 \((\mathfrak{g},\mathfrak{h})\) 的一个抛物子代数，Q 是 \(\alpha \in \mathrm{P}\) 中使 \(-\alpha \notin \mathrm{P}\) 者组成的集合，并设 \(\mathfrak{s} = \mathfrak{h} + \mathfrak{g}^{\mathrm{P}\cap (-\mathrm{P})}\) 。则 \(\mathfrak{p} = \mathfrak{s} \oplus \mathfrak{g}^{\mathrm{Q}}\) ，\(\mathfrak{s}\) 在 \(\mathfrak{g}\) 中是约化的，且 \(\mathfrak{g}^{\mathrm{Q}}\) 是 \(\mathfrak{p}\) 的最大幂零理想并且是 \(\mathfrak{p}\) 的幂零根基。\(\mathfrak{p}\) 的中心为零。

由命题 2，\(\mathfrak{s}\) 在 \(\mathfrak{g}\) 中是约化的，且 \(\mathfrak{g}^{\mathrm{Q}}\) 是 \(\mathfrak{p}\) 的一个幂零理想。若 \(\mathfrak{n}\) 是 \(\mathfrak{p}\) 的最大幂零理想，则 \(\mathfrak{g}^{\mathrm{Q}} \subset \mathfrak{n} \subset \mathfrak{h} + \mathfrak{g}^{\mathrm{Q}}\) （命题 2 (i)）；若 \(x \in \mathfrak{n} \cap \mathfrak{h}\) ，则 \(\operatorname{ad}_{\mathfrak{p}} x\) 是幂零的，故 \(\alpha(x) = 0\) 对一切 \(\alpha \in \mathrm{P}\) 成立，从而 \(x = 0\) ；这就证明了 \(\mathfrak{n} = \mathfrak{g}^{\mathrm{Q}}\) 。由于 \([\mathfrak{h}, \mathfrak{g}^{\mathrm{Q}}] = \mathfrak{g}^{\mathrm{Q}}\) ，\(\mathfrak{p}\) 的幂零根基包含 \(\mathfrak{g}^{\mathrm{Q}}\) ，从而等于 \(\mathfrak{g}^{\mathrm{Q}}\) 。设 \(z = h + \sum_{\alpha \in \mathrm{P}} u_{\alpha}\) （其中 \(h \in \mathfrak{h}, u_{\alpha} \in \mathfrak{g}^{\alpha}\) ）是 \(\mathfrak{p}\) 的中心的一个元素。对一切 \(h' \in \mathfrak{h}\) ，有 \(0 = [h', z] = \sum \alpha(h')u_{\alpha}\) ，故 \(u_{\alpha} = 0\) 对一切 \(\alpha \in \mathrm{P}\) 成立；由此推出 \([h, \mathfrak{g}^{\beta}] = 0\) 对一切 \(\beta \in \mathrm{P}\) 成立，故 \(h = 0\) 。

\subsection{非分裂情形}

命题 14. 设 \(k'\) 是 \(k\) 的扩张，\(\mathfrak{g}' = \mathfrak{g} \otimes_k k'\) 。设 \(\mathfrak{m}\) 是 \(\mathfrak{g}\) 的子代数，\(\mathfrak{m}' = \mathfrak{m} \otimes_k k'\) 。若 \(\mathfrak{m}'\) 是 \(\mathfrak{g}'\) 的抛物（相应地 Borel）子代数，则 \(\mathfrak{m}\) 是 \(\mathfrak{g}\) 的抛物（相应地 Borel）子代数。

由命题 8 与 11，只需证明 m 包含 g 的一个分裂嘉当子代数即可。设 b 是 g 的一个 Borel 子代数。则 \(b' = b \otimes_{k} k'\) 是 \(\mathfrak{g}'\) 的 Borel 子代数，故 \(\mathfrak{m}' \cap \mathfrak{b}'\) 包含 \(\mathfrak{g}'\) 的一个嘉当子代数（命题 10）。设 \(\mathfrak{t}\) 是 \(\mathfrak{m} \cap \mathfrak{b}\) 的一个嘉当子代数。则 \(\mathfrak{t} \otimes_k k'\) 是 \(\mathfrak{m}' \cap \mathfrak{b}'\) 的、从而是 \(\mathfrak{g}'\) 的嘉当子代数（第七章 §3 第 3 节 命题 3）。于是，\(\mathfrak{t}\) 是 \(\mathfrak{g}\) 的嘉当子代数，且由于它含于 \(\mathfrak{b}\) ，它是分裂的。

定义 3. 设 \(\mathfrak{a}\) 是半单（或更一般地，约化）李代数，\(\bar{k}\) 是 \(k\) 的一个代数闭包。子代数 \(\mathfrak{m}\) （\(\mathfrak{a}\) 的）称为抛物的（相应地 Borel 的），如果 \(\mathfrak{m} \otimes_k \bar{k}\) 是 \(\mathfrak{a} \otimes_k \bar{k}\) 的抛物（相应地 Borel）子代数。

若 \(\mathfrak{a}\) 是可分裂的，则命题 14 表明这个定义与定义 2（相应地与定义 1）等价。

命题 15. 设 \(\mathfrak{a}\) 是约化李代数，\(k'\) 是 \(k\) 的扩张，\(\mathfrak{m}\) 是 \(\mathfrak{a}\) 的子代数。则 \(\mathfrak{m}\) 是 \(\mathfrak{a}\) 的抛物（相应地 Borel）子代数当且仅当 \(\mathfrak{m} \otimes_k k'\) 是 \(\mathfrak{a} \otimes_k k'\) 的抛物（相应地 Borel）子代数。

这由命题 14 立即得出。

\section{由既约根系定义的分裂半单李代数}

\subsection{带标架的半单李代数}

命题 1. 设 \((\mathfrak{g},\mathfrak{h})\) 是分裂半单李代数，R 是它的根系，B 是 R 的基，\((n(\alpha ,\beta))_{\alpha ,\beta \in \mathrm{B}}\) 是对应的嘉当矩阵。对一切 \(\alpha \in \mathrm{B}\) ，设 \(X_{\alpha}\in \mathfrak{g}^{\alpha}\) ，\(X_{-\alpha}\in \mathfrak{g}^{-\alpha}\) 。则对 \(\alpha ,\beta \in \mathrm{B}\) ，

\[
[ H _ {\alpha}, H _ {\beta} ] = 0\tag{1}
\]

\[
[ H _ {\alpha}, X _ {\beta} ] = n (\beta , \alpha) X _ {\beta}\tag{2}
\]

\[
[ H _ {\alpha}, X _ {- \beta} ] = - n (\beta , \alpha) X _ {- \beta}\tag{3}
\]

\[
[ X _ {- \alpha}, X _ {\beta} ] = 0 \quad \text { if } \alpha \neq \beta\tag{4}
\]

\[
(\operatorname{ad} X _ {\alpha}) ^ {1 - n (\beta , \alpha)} X _ {\beta} = 0 \quad \text { if } \alpha \neq \beta\tag{5}
\]

\[
(\operatorname{ad} X _ {- \alpha}) ^ {1 - n (\beta , \alpha)} X _ {- \beta} = 0 \quad \text { if } \alpha \neq \beta .\tag{6}
\]

族 \((H_{\alpha})_{\alpha \in \mathrm{B}}\) 是 \(\mathfrak{h}\) 的一个基。若 \(X_{\alpha} \neq 0\) 且 \(X_{-\alpha} \neq 0\) 对一切 \(\alpha \in \mathrm{B}\) 成立，则李代数 \(\mathfrak{g}\) 由诸 \(X_{\alpha}\) 与诸 \(X_{-\alpha}\) （ \(\alpha \in \mathrm{B}\) ）生成。

（回忆，若 \(\alpha, \beta \in \mathrm{B}\) 且 \(\alpha \neq \beta\) ，则 \(n(\beta, \alpha)\) 是一个 \(\leq 0\) 的整数，故公式 (5) 与 (6) 有意义。）

公式 (1)、(2) 与 (3) 是明显的。若 \(\alpha \neq \beta\) ，则 \(\beta - \alpha\) 不是根，因为 \(R\) 的每个元素都是 \(B\) 的元素的具整数系数的线性组合，且这些系数全同号（第六章 §1 第 6 节 定理 3）。这就证明了 (4)。鉴于第六章 §1 第 3 节 命题 9，这还证明了 \(\alpha\) -链（由 \(\beta\) 定义的）是

\[
\{\beta , \beta + \alpha , \dots , \beta - n (\beta , \alpha) \alpha \};
\]

从而 \(\beta + (1 - n(\beta, \alpha))\alpha \notin \mathrm{R}\) ，这就证明了 (5)。等式 (6) 用类似的方式建立。族 \((H_{\alpha})_{\alpha \in \mathrm{B}}\) 是 \(\mathrm{R}^{\vee}\) 的、从而是 \(\mathfrak{h}\) 的基。若 \(X_{\alpha} \neq 0\) 且 \(X_{-\alpha} \neq 0\) 对一切 \(\alpha \in \mathrm{B}\) 成立，则 \([X_{\alpha}, X_{-\alpha}] = \lambda_{\alpha}H_{\alpha}\) ，其中 \(\lambda_{\alpha} \neq 0\) ，故最后的断言由 §3 第 3 节 命题 9 (iii) 得出。

定义 1. 设 \((\mathfrak{g},\mathfrak{h})\) 是分裂半单李代数，R 是它的根系。\((\mathfrak{g},\mathfrak{h})\) 的一个标架是指一个偶 \((\mathrm{B},(X_{\alpha})_{\alpha \in \mathrm{B}})\) ，其中 B 是 R 的一个基，并且对一切 \(\alpha \in \mathrm{B}\) ，\(X_{\alpha}\) 是 \(\mathfrak{g}^{\alpha}\) 的一个非零元素。带标架的半单李代数是指一个序列 \((\mathfrak{g},\mathfrak{h},\mathrm{B},(X_{\alpha})_{\alpha \in \mathrm{B}})\) ，其中 \((\mathfrak{g},\mathfrak{h})\) 是分裂半单李代数，且 \((\mathrm{B},(X_{\alpha})_{\alpha \in \mathrm{B}})\) 是 \((\mathfrak{g},\mathfrak{h})\) 的一个标架。

\(\mathfrak{g}\) 的一个标架是指 \((\mathfrak{g},\mathfrak{h})\) 的一个标架，其中 \(\mathfrak{h}\) 是 \(\mathfrak{g}\) 的一个分裂嘉当子代数。

设 \(a_1 = (\mathfrak{g}_1, \mathfrak{h}_1, \mathrm{B}_1, (X_\alpha^1)_{\alpha \in \mathrm{B}_1})\) 与 \(a_2 = (\mathfrak{g}_2, \mathfrak{h}_2, \mathrm{B}_2, (X_\alpha^2)_{\alpha \in \mathrm{B}_2})\) 是带标架的半单李代数。从 \(a_1\) 到 \(a_2\) 的一个同构是指一个同构 \(\varphi\) （从 \(\mathfrak{g}_1\) 到 \(\mathfrak{g}_2\) ），它把 \(\mathfrak{h}_1\) 映到 \(\mathfrak{h}_2\) ，把 \(\mathrm{B}_1\) 映到 \(\mathrm{B}_2\) ，并且把 \(X_\alpha^1\) 映到 \(X_{\psi \alpha}^2\) 对一切 \(\alpha \in \mathrm{B}_1\) 成立（其中 \(\psi\) 是 \(\varphi | \mathfrak{h}_1\) 的反步映射）。这时，称 \(\varphi\) 把标架 \((\mathrm{B}_1, (X_\alpha^1)_{\alpha \in \mathrm{B}_1})\) 变换为标架 \((\mathrm{B}_2, (X_\alpha^2)_{\alpha \in \mathrm{B}_2})\) 。

若 \((\mathrm{B},(X_{\alpha})_{\alpha\in\mathrm{B}})\) 是 \((\mathfrak{g},\mathfrak{h})\) 的一个标架，则对一切 \(\alpha\in B\) ，存在唯一的元素 \(X_{-\alpha}\) （\(g^{-\alpha}\) 中的）使 \([X_{\alpha},X_{-\alpha}]=-H_{\alpha}\) （§2 第 2 节 定理 1 (iv)）。族 \((X_{\alpha})_{\alpha\in\mathrm{B}\cup(-\mathrm{B})}\) 称为由该标架定义的生成族（参见命题 1）。它也是由标架 \((-\mathrm{B},(X_{\alpha})_{\alpha\in-\mathrm{B}})\) 定义的生成族。对一切 \(\alpha\in\mathrm{B}\cup(-\mathrm{B})\) ，设 \(t_{\alpha}\in k^{*}\) ，并设 \(t_{\alpha}t_{-\alpha}=1\) 对一切 \(\alpha\in B\) 成立。则 \((t_{\alpha}X_{\alpha})_{\alpha\in\mathrm{B}\cup(-\mathrm{B})}\) 是由标架 \((\mathrm{B},(t_{\alpha}X_{\alpha})_{\alpha\in\mathrm{B}})\) 定义的生成族。

\subsection{一个预备性构造}

在本节及下一节中，我们用 R 表示向量空间 V 中的一个既约根系，用 B 表示 R 的一个基。我们用 \((n(\alpha,\beta))_{\alpha,\beta\in\mathbb{B}}\) 表示相对于 B 的嘉当矩阵。回忆 \(n(\alpha,\beta)=\langle\alpha,\beta^{\vee}\rangle\) 。我们将证明 R 是某个分裂半单李代数的根系，且该李代数在同构意义下唯一。我们主要将考虑由命题 1 中的关系所定义的李代数。

本节的构造适用于 k 上任何具非零行列式的方阵 \((n(\alpha,\beta))_{\alpha,\beta\in\mathbb{B}}\) ，只要 \(n(\alpha,\alpha)=2\) 对一切 \(\alpha\in B\) 成立（参见第六章 §1 第 10 节 公式 (14)）。

设 E 是集合 B 在 k 上的自由结合代数。回忆 E 是 N-分次的（《代数》第三章 §3 第 1 节 例 3）。我们将对每个 \(\alpha\in B\) 联系向量空间 E 的自同态 \(X_{-\alpha}^{0}, H_{\alpha}^{0}, X_{\alpha}^{0}\) ，其次数分别为 1、0、-1。对 B 的元素组成的任何字 \((\alpha_{1},\ldots,\alpha_{n})\) ，置

\[
X _ {- \alpha} ^ {0} (\alpha_ {1}, \ldots , \alpha_ {n}) = (\alpha , \alpha_ {1}, \ldots , \alpha_ {n})\tag{7}
\]

\[
H _ {\alpha} ^ {0} \left(\alpha_ {1}, \dots , \alpha_ {n}\right) = \left(- \sum_ {i = 1} ^ {n} n \left(\alpha_ {i}, \alpha\right)\right) \left(\alpha_ {1}, \dots , \alpha_ {n}\right).\tag{8}
\]

另一方面，\(X_{\alpha}^{0}(\alpha_{1},\ldots ,\alpha_{n})\) 借助下列公式对 \(n\) 用归纳法来定义

\[
X _ {\alpha} ^ {0} (\alpha_ {1}, \ldots , \alpha_ {n}) = (X _ {- \alpha_ {1}} ^ {0} X _ {\alpha} ^ {0} - \delta_ {\alpha , \alpha_ {1}} H _ {\alpha} ^ {0}) (\alpha_ {2}, \ldots , \alpha_ {n})\tag{9}
\]

其中 \(\delta_{\alpha, \alpha_1}\) 是 Kronecker 符号；约定 \(X_{\alpha}^{0}(\alpha_{1}, \ldots, \alpha_{n})\) 在 \((\alpha_{1}, \ldots, \alpha_{n})\) 为空字时为零。

引理 1. 对一切 \(\alpha, \beta \in \mathbf{B}\) ，我们有

\[
[ X _ {\alpha} ^ {0}, X _ {- \alpha} ^ {0} ] = - H _ {\alpha} ^ {0}\tag{10}
\]

\[
[ H _ {\alpha} ^ {0}, H _ {\beta} ^ {0} ] = 0\tag{11}
\]

\[
[ H _ {\alpha} ^ {0}, X _ {\beta} ^ {0} ] = n (\beta , \alpha) X _ {\beta} ^ {0}\tag{12}
\]

\[
[ H _ {\alpha} ^ {0}, X _ {- \beta} ^ {0} ] = - n (\beta , \alpha) X _ {- \beta} ^ {0}\tag{13}
\]

\[
[ X _ {\alpha} ^ {0}, X _ {- \beta} ^ {0} ] = 0 \text {   if   } \alpha \neq \beta .\tag{14}
\]

事实上，关系式 (9) 可以写成

\[
(X _ {\alpha} ^ {0} X _ {- \alpha_ {1}} ^ {0}) (\alpha_ {2}, \ldots , \alpha_ {n}) = (X _ {- \alpha_ {1}} ^ {0} X _ {\alpha} ^ {0}) (\alpha_ {2}, \ldots , \alpha_ {n}) - \delta_ {\alpha , \alpha_ {1}} H _ {\alpha} ^ {0} (\alpha_ {2}, \ldots , \alpha_ {n})
\]

这证明了 (10) 与 (14)。关系式 (11) 是明显的。其次

\[
\begin{array}{r l} [ H _ {\alpha} ^ {0}, X _ {- \beta} ^ {0} ] (\alpha_ {1}, \ldots , \alpha_ {n}) & = H _ {\alpha} ^ {0} (\beta , \alpha_ {1}, \ldots , \alpha_ {n}) + \left(\sum_ {i = 1} ^ {n} n (\alpha_ {i}, \alpha)\right) (\beta , \alpha_ {1}, \ldots , \alpha_ {n}) \\ & = - n (\beta , \alpha) (\beta , \alpha_ {1}, \ldots , \alpha_ {n}) \\ & = - n (\beta , \alpha) X _ {- \beta} ^ {0} (\alpha_ {1}, \ldots , \alpha_ {n}) \end{array}
\]

从而得 (13)。最后，

\[
\begin{array}{l} 0 = [ H _ {\alpha} ^ {0}, [ X _ {\beta} ^ {0}, X _ {- \gamma} ^ {0} ] ] \quad \text { by (10), (11), (14)} \\ \qquad = [ [ H _ {\alpha} ^ {0}, X _ {\beta} ^ {0} ], X _ {- \gamma} ^ {0} ] + [ X _ {\beta} ^ {0}, [ H _ {\alpha} ^ {0}, X _ {- \gamma} ^ {0} ] ] \\ \qquad = [ [ H _ {\alpha} ^ {0}, X _ {\beta} ^ {0} ] - n (\gamma , \alpha) X _ {\beta} ^ {0}, X _ {- \gamma} ^ {0} ] \quad \text { by (13)} \\ \qquad = [ [ H _ {\alpha} ^ {0}, X _ {\beta} ^ {0} ] - n (\beta , \alpha) X _ {\beta} ^ {0}, X _ {- \gamma} ^ {0} ] \quad \text { by (14); } \end{array}\tag{15}
\]

而对空字加以考虑立即给出

\[
([ H _ {\alpha} ^ {0}, X _ {\beta} ^ {0} ] - n (\beta , \alpha) X _ {\beta} ^ {0}) (\emptyset) = 0
\]

故 (15) 蕴含

\[
\left(\left[ H _ {\alpha} ^ {0}, X _ {\beta} ^ {0} \right] - n (\beta , \alpha) X _ {\beta} ^ {0}\right) X _ {- \gamma_ {1}} ^ {0} X _ {- \gamma_ {2}} ^ {0} \dots X _ {- \gamma_ {n}} ^ {0} (\varnothing) = 0
\]

对一切 \(\gamma_1, \ldots, \gamma_n \in \mathrm{B}\) 成立；这就证明了 (12)。

引理 2. 自同态 \(X_{\alpha}^{0}\) ，\(H_{\beta}^{0}\) ，\(X_{-\gamma}^{0}\) （其中 \(\alpha, \beta, \gamma \in \mathrm{B}\) ）线性无关。

由于 \(X_{-\alpha}^{0}(\varnothing) = \alpha\) ，显然诸 \(X_{-\alpha}^{0}\) 线性无关。设 \(\sum_{\alpha} a_{\alpha} H_{\alpha}^{0} = 0\) ；则对一切 \(\beta \in B\) ，

\[
0 = \left[ \sum_ {\alpha} a _ {\alpha} H _ {\alpha} ^ {0}, X _ {- \beta} ^ {0} \right] = - \sum_ {\alpha} a _ {\alpha} n (\beta , \alpha) X _ {- \beta} ^ {0};
\]

由于 \(\operatorname{det}(n(\beta, \alpha)) \neq 0\) ，推出 \(a_{\alpha} = 0\) 对一切 \(\alpha\) 成立。设 \(\sum_{\alpha} a_{\alpha} X_{\alpha}^{0} = 0\) 。鉴于公式 (7)、(8)、(9)，

\[
X _ {\alpha} ^ {0} (\beta) = 0,
\]

\[
X _ {\alpha} ^ {0} (\beta , \beta) = 2 \delta_ {\alpha \beta} \beta
\]

对一切 \(\beta \in \mathrm{B}\) 成立。由此推出 \(a_{\beta} = 0\) 对一切 \(\beta\) 成立。由于 \(X_{\alpha}^{0}\) ，\(H_{\alpha}^{0}\) ，\(X_{-\alpha}^{0}\) 的次数分别为 \(-1, 0, 1\) ，引理由前面所述即得。

设 I 是集合 \(B \times \{-1,0,1\}\) 。置 \(x_{\alpha} = (\alpha, -1)\) ，\(h_{\alpha} = (\alpha, 0)\) ，以及 \(x_{-\alpha} = (\alpha, 1)\) 。设 a 是由生成族 I 与下列关系元集合 R 所定义的李代数：

\[
\begin{array}{l} \left[ h _ {\alpha}, h _ {\beta} \right] \\ \left[ h _ {\alpha}, x _ {\beta} \right] - n (\beta , \alpha) x _ {\beta} \\ \left[ h _ {\alpha}, x _ {- \beta} \right] + n (\beta , \alpha) x _ {- \beta} \\ \left[ x _ {\alpha}, x _ {- \alpha} \right] + h _ {\alpha} \\ \left[ x _ {\alpha}, x _ {- \beta} \right] \quad \text {if} \alpha \neq \beta \end{array}
\]

（参见第二章 §2 第 3 节）。由引理 1，存在唯一的线性表示 \(\rho\) （\(\mathfrak{a}\) 在 \(\mathrm{E}\) 上的），使得

\[
\rho (x _ {\alpha}) = X _ {\alpha} ^ {0}, \rho (h _ {\alpha}) = H _ {\alpha} ^ {0}, \rho (x _ {- \alpha}) = X _ {- \alpha} ^ {0}.
\]

鉴于引理 2，这就证明了下列结果：

引理 3. \(\mathfrak{a}\) 中元素 \(x_{\alpha}, h_{\beta}, x_{-\gamma}\) （其中 \(\alpha, \beta, \gamma \in \mathrm{B}\) ）的典则像线性无关。

在下文中，我们把 \(x_{\alpha}\) ，\(h_{\alpha}\) ，\(x_{-\alpha}\) 与它们在 a 中的典则像等同起来。

引理 4. 存在唯一的对合自同构 \(\theta\) （\(\mathfrak{a}\) 的），使得

\[
\theta (x _ {\alpha}) = x _ {- \alpha}, \theta (x _ {- \alpha}) = x _ {\alpha}, \theta (h _ {\alpha}) = - h _ {\alpha}
\]

对一切 \(\alpha\in B\) 成立。

事实上，自由李代数 \(L(I)\) 存在一个满足这些条件的对合自同构。它使 \(R \cup (-R)\) 稳定，因而通过过渡到商定义了 a 的一个满足引理条件的对合自同构。唯一性由 a 由元素 \(x_{\alpha}\) ，\(h_{\alpha}\) ，\(x_{-\alpha}\) （ \(\alpha \in B\) ）生成这一事实得出。

这个自同构称为 a 的典则对合自同构。

设 Q 是 R 的根权组成的集合；它是一个以 B 为基的自由 Z-模（第六章 §1 第 9 节）。在自由李代数 L(I) 上存在一个 Q 型分次，使 \(x_{\alpha}\) ，\(h_{\alpha}\) ，\(x_{-\alpha}\) 的次数分别为 \(\alpha\) ，0，- \(\alpha\) （第二章 §2 第 6 节）。而 R 的元素都是齐次的。因此，在 a 上存在唯一的、与 a 的李代数结构相容的 Q 型分次，使 \(x_{\alpha}\) ，\(h_{\alpha}\) ，\(x_{-\alpha}\) 的次数分别为 \(\alpha\) ，0，- \(\alpha\) 。对任何 \(\mu \in Q\) ，用 \(a^{\mu}\) 表示 a 中次数为 \(\mu\) 的齐次元素组成的集合。

引理 5. 设 \(z \in \mathfrak{a}\) 。则 \(z \in \mathfrak{a}^{\mu}\) 当且仅当 \([h_{\alpha}, z] = \langle \mu, \alpha^{\vee} \rangle z\) 对一切 \(\alpha \in \mathrm{B}\) 成立。

对 \(\mu \in \mathrm{Q}\) ，设 \(\mathfrak{a}^{(\mu)}\) 是 \(x \in \mathfrak{a}\) 中使 \([h_{\alpha}, x] = \langle \mu, \alpha^{\vee} \rangle x\) 对一切 \(\alpha \in \mathrm{B}\) 成立者组成的集合。诸 \(\mathfrak{a}^{(\mu)}\) 之和是直和。因此，为证明引理，只需证明 \(\mathfrak{a}^{\mu} \subset \mathfrak{a}^{(\mu)}\) 。设 \(\alpha \in \mathrm{B}\) 。自同态 \(u\) （向量空间 \(\mathfrak{a}\) 的，它使 \(u|\mathfrak{a}^{\mu} = \langle \mu, \alpha^{\vee} \rangle\) .1）是 \(\mathfrak{a}\) 的一个导子，它使 \(ux = (\operatorname{ad} h_{\alpha})\) . \(x\) 对 \(x = x_{\beta}\) ，\(x = h_{\beta}\) ，\(x = x_{-\beta}\) 成立；故 \(u = \operatorname{ad} h_{\alpha}\) ，这就证明了我们的断言。

注. 由引理 5 推出 \(\mathfrak{a}\) 的每个理想都是齐次的，因为它在 ad \(h_{\alpha}\) 之下稳定。

用 \(Q_{+}\) （相应地 \(Q_{-}\) ）表示 B 中元素的具正（相应地负）整系数、且不全为零的线性组合组成的集合。置 \(a_{+} = \sum_{\mu \in Q_{+}} a^{\mu}\) 与 \(a_{-} = \sum_{\mu \in Q_{-}} a^{\mu}\) 。由于 \(Q_{+} + Q_{+} \subset Q_{+}\) 且 \(Q_{-} + Q_{-} \subset Q_{-}\) ，\(a_{+}\) 与 \(a_{-}\) 是 a 的李子代数。

命题 2. (i) 李代数 \(\mathfrak{a}_{+}\) 由族 \((x_{\alpha})_{\alpha \in \mathbb{B}}\) 生成。

\begin{enumerate}
\def\labelenumi{(\roman{enumi})}
\setcounter{enumi}{1}
\item
  李代数 \(\mathfrak{a}_{-}\) 由族 \((x_{-\alpha})_{\alpha \in \mathrm{B}}\) 生成。
\item
  族 \((h_{\alpha})_{\alpha \in \mathrm{B}}\) 是向量空间 \(\mathfrak{a}^0\) 的一个基。
\item
  向量空间 \(\mathfrak{a}\) 是 \(\mathfrak{a}_{+},\mathfrak{a}^{0},\mathfrak{a}_{-}\) 的直和。
\end{enumerate}

设 \(\mathfrak{r}\) （相应地 \(\mathfrak{n}\) ）是 \(\mathfrak{a}\) 的由 \((x_{\alpha})_{\alpha \in \mathrm{B}}\) （相应地 \((x_{-\alpha})_{\alpha \in \mathrm{B}}\) ）生成的李子代数，并设 \(\mathfrak{h}\) 是 \(\mathfrak{a}\) 的由 \((h_{\alpha})_{\alpha \in \mathrm{B}}\) 生成的向量子空间。由于诸 \(x_{\alpha}\) 是 \(\mathfrak{a}_{+}\) 的齐次元素，\(\mathfrak{r}\) 是 \(\mathfrak{a}_{+}\) 的一个分次子代数；于是 \([\mathfrak{h},\mathfrak{r}]\subset \mathfrak{r}\) ，故 \(\mathfrak{h} + \mathfrak{r}\) 是 \(\mathfrak{a}\) 的一个子代数；由于

\[
[ x _ {- \alpha}, x _ {\beta} ] = \delta_ {\alpha \beta} h _ {\alpha},
\]

\([x_{-\alpha}, \mathfrak{r}] \subset \mathfrak{h} + \mathfrak{r}\) 对一切 \(\alpha \in B\) 成立。类似地，\(\mathfrak{n}\) 是 \(\mathfrak{a}_{-}\) 的一个分次子代数，有 \([\mathfrak{h}, \mathfrak{n}] \subset \mathfrak{n}\) ，\(\mathfrak{h} + \mathfrak{n}\) 是 \(\mathfrak{n}\) 的一个子代数，且 \([x_{\alpha}, \mathfrak{n}] \subset \mathfrak{h} + \mathfrak{n}\) 对一切 \(\alpha \in B\) 成立。置 \(a' = r + h + n\) 。前面的论证表明，\(a'\) 在 ad \(x_{\alpha}\) 、ad \(h_{\alpha}\) 与 ad \(x_{-\alpha}\) 之下对一切 \(\alpha \in B\) 稳定，从而是 a 的一个理想。由于 \(a'\) 包含 \(x_{\alpha}\) ，\(h_{\alpha}\) ，\(x_{-\alpha}\) 对一切 \(\alpha \in B\) 成立，故 \(a' = a\) 。由此推出包含关系 \(r \subset a_{+}\) ，\(h \subset a^{0}\) ，\(n \subset a_{-}\) 都是等式，这就证明了命题。

命题 3. 李代数 \(\mathfrak{a}_{+}\) （相应地 \(\mathfrak{a}_{-}\) ）是以 \((x_{\alpha})_{\alpha \in \mathrm{B}}\) （相应地 \((x_{-\alpha})_{\alpha \in \mathrm{B}}\) ）为基础族的自由李代数（参见第二章 §2 第 3 节）。

设 L 是 E 的由 B 生成的李子代数。由第二章 §3 定理 1，L 可与由 B 生成的自由李代数等同。E 在其自身上的左正则表示显然是单射的，且通过限制到 L 上定义了李代数 L 在 E 上的一个单射表示 \(\rho'\) 。设 \(\varphi\) 是从 L 到 \(\mathfrak{a}_{-}\) 的唯一的同态，它把 \(\alpha\) 映为 \(x_{-\alpha}\) 对一切 \(\alpha \in B\) 成立。则对一切 \(\alpha \in B\) ，\(\rho(\varphi(\alpha))\) 是 E 上用 \(\alpha\) 作左乘的自同态，故 \(\rho \circ \varphi = \rho'\) ，这证明 \(\varphi\) 是单射的。这样，\((x_{-\alpha})_{\alpha \in B}\) 是 \(\mathfrak{a}_{-}\) 的一个基础族。由于 \(\theta(x_{-\alpha}) = x_{\alpha}\) 对一切 \(\alpha\) 成立（参见引理 4），\((x_{\alpha})_{\alpha \in B}\) 是 \(\mathfrak{a}_{+}\) 的一个基础族。